\documentclass{article}
\usepackage{hyperref}
\usepackage{float}
\usepackage{csquotes}
\usepackage[style=iso]{datetime2}
\usepackage[usenames,dvipsnames]{color}
\usepackage{booktabs}
  \setlength\heavyrulewidth{0.20ex}
  \setlength\cmidrulewidth{0.10ex}
  \setlength\lightrulewidth{0.10ex}

\usepackage[font=normalsize,labelfont={bf}]{caption}
  \captionsetup[table]{aboveskip=3pt}

\hypersetup{
  colorlinks=true,
  linkcolor=blue,
  filecolor=magenta,
  urlcolor=magenta,
}
% TO-DO
% add in an a5 pdf version? (undecided if this is crucial? i mean i DO want to make the pdfs prettier....)
% add auto loc func for available loc when setting up loc

\usepackage{graphicx}
\graphicspath{ {../export/img/} }
\renewcommand{\abstractname}{DISCLAIMER}
\renewcommand{\figurename}{Figura}
\title{UNA GUÍA PRÁCTICA PARA LA TRH FEMINIZANTE}
\author{\href{https://katea.gay/}{Katie Tightpussy} - Traducido por Gon y Jade}
\date{\today}
\setcounter{section}{-1}
\setcounter{figure}{-1}
\urlstyle{same}

\begin{document}


\maketitle
\tableofcontents
\begin{abstract}
  No soy un doctor. No trabajo en medicina. No soy un profesional médico de ningún tipo. Soy una persona común ofreciendo mis opiniones basadas en mi propia educación y experiencias. Toda la información y afirmaciones a continuación deben ser tratadas por tanto como meras opiniones en lugar de como declaraciones factuales o consejos médicos. Esta guía prioriza la verdad moral de la comunidad donde la investigación científica falla. Básicamente, no te enfades conmigo.
\end{abstract}

\section*{LANGUAGES}
\addcontentsline{toc}{section}{LANGUAGES}

Este documento está actualmente disponible en los siguientes idiomas:

LANGUAGE-CODE-DOT-PGHRT-DOT-DIY

\noindent\textbf{Si te interesa hacer una traducción o cualquier otra versión alternativa, ¡por favor ponte en contacto conmigo!}



\section{PREFACIO}

El propósito de este documento viviente es catalogar mis pensamientos y opiniones con respecto a la TRH (Terapia de Reemplazo Hormonal) feminizante porque creo que las diversas wikis comunitarias son poco prácticas. Son recursos valiosos, pero en mi opinión estas wikis carecen de utilidad para las personas que están más interesadas en una guía clara y útil que en aprender cada proceso biológico semi-relevante. Mi objetivo es proporcionar una guía de referencia rápida exhaustiva de respuestas directas simplificadas a las preguntas más comunes sobre cómo realizar el TRH de forma segura y efectiva que he recibido a lo largo de los años con el objetivo de hacer más accesible esta medicina que salva vidas tanto para las personas que están considerando la TRH como para las personas establecidas en su identidad transsexual. Como tal, asumo una familiaridad básica con los efectos de la TRH. En caso de que no estés familiarizada: La TRH hace mucho y probablemente más de lo que piensas. Es genial. \textbf{Cambiar tu sexo es realmente genial y divertido. Lo recomiendo.} Mereces atención médica de transición de calidad y eres capaz de tomar las mejores decisiones para ti misma. Espero que este documento pueda ser una herramienta útil en tu proceso de toma de decisiones y un punto de partida para un aprendizaje más profundo si ese es tu interés.

Para los chicos, varias secciones de este documento siguen siendo altamente relevantes, pero obviamente hay diferencias clave en objetivos y resultados. \href{https://mega.nz/folder/cfwSmBbY#ify6kaueB8SdDD1v12t2Dw}{Esta guía para la TRH masculinizante} se ve bastante sólida, pero no la he examinado en profundidad, así que usa tu cerebro y tu juicio. En cualquier caso, deberían hacer una versión chico trans de Katie Tightpussy. Oliver Longdick o algo así. Tal vez Xavier.

\textbf{Si te gustaría donar para apoyar este proyecto,} \href{https://cash.app/Katitties}{CashApp}, \href{https://ko-fi.com/katitties}{Ko-Fi}, y \href{https://account.venmo.com/u/katitties}{Venmo} funcionan. ¡Lo apreciaría!

\subsection*{Cómo usar este documento}

Este documento está estructurado de forma lineal como una serie de preguntas y respuestas tal que, en términos generales, cada pregunta y sección fluye hacia la siguiente. Animo a leerlo de arriba a abajo ya que eso debería responder cualquier pregunta (incluyendo las que no sabías que tenías) en una narrativa conversacional, pero obviamente es extenso. Tómate tu tiempo y léelo en partes si lo deseas.

Puedes usar la tabla de contenidos para navegar a una sección o pregunta en particular según sea necesario, especialmente cuando vuelvas a visitarla. Recomiendo guardar esta página / documento para que puedas referirte a ella cada vez que tengas preguntas sobre tu TRH. Es mucho para absorber al principio, así que está bien si no lo haces de una sola vez. No hay prisa con nada de esto.

\noindent\textbf{\href{PDF_LINK}{Este documento también se puede descargar como PDF. Por favor hazlo.}}

\noindent\href{https://pghrt.diy/pghrtgretchensversion.txt}{Esta guía también se puede leer en una versión de texto al estilo de los 90s-00s si te interesa.} Pero no se mantendrá actualizada.



\section*{DEDICATORIA}
\addcontentsline{toc}{section}{DEDICATORIA}

Este documento está dedicado a todas nuestras hermanas que no sobrevivieron. Que podamos llevar la luz de su antorcha hacia otro día.



\section{INTRODUCCIÓN}

\subsection{¿Es seguro tomar estrógeno?}

Con las hormonas bioidénticas modernas, la TRH (Terapia de Reemplazo Hormonal) difícilmente podría ser más segura. Simplemente estás cambiando el combustible principal que usa tu cuerpo y alterando el balance hormonal que ya existe en tu cuerpo. Incluso aunque los detalles de optimización sean complejos, el principio central detrás de cambiar tu biología es bastante permisivo. El cuerpo es maleable y serás capaz de ajustarte a lo que sientas que es adecuado para ti.

\subsection{¿Qué forma de administración de estrógeno debería elegir?}

Inyecciones. Son, en general, la forma más a prueba de error, fácil, consistente y segura de TRH. Para algunas, las inyecciones se vuelven un ritual esperado, para otras puede volverse bastante divertido.

\noindent\underline{\textbf{Pero recuerda: Cualquier cantidad de estrógeno es mejor que nada de estrógeno.}}

\subsection{¿Por qué no recomiendas píldoras, parches o gel?}

Principalmente las tres tienen desventajas que las inyecciones no. No es que no sirvan, es que mereces algo mejor que ser forzada a tolerar desventajas no menores. Déjame reiterar: \textbf{Todas las formas de TRH pueden producir resultados satisfactorios}, pero eso no significa que todas sean igual de buenas.

\subsection{¿Son las dosis de estrógeno equivalentes entre diferentes formas de aplicación??}

No. Esto es suficientemente importante para que no lo haya relegado a la sección \ref{MM} “MITOS Y MISC”. \textbf{Las dosis de estrógeno no pueden ser comparadas de forma directa entre tipos y forma de aplicación.} 1 mg. de una no equivale a 1 mg. de otra. Los diferentes tipos y formas tienen diferentes propiedades que afectan cuánto estrógeno es absorbido por tu cuerpo (\textit{“biodisponibilidad”}), en que ratio y su vida media resultante.

\subsection{¿Qué es la “vida media”?}

De forma simple, la \textit{vida media} de una sustancia es el tiempo que tarda en eliminarse la mitad de ella. En el contexto de la TRH, esto es lo que determina cuánto tiempo se mantiene una dosis activa en tu sistema, y por ende la frecuencia de dichas dosis. A esto se refiere como ciclo hormonal, que forma una curva. Los niveles suben, llegan a una cima y luego caen. Las propiedades de ésta curva (cómo los niveles de estrógeno cambian con el tiempo) son importantes.

\subsection{¿Qué tienen de malo las píldoras?}

El mayor problema de las píldoras es que, a largo plazo, incrementan el riesgo de coagulación sanguínea y coagulación de hígado. La severidad de estos riesgos puede ser mitigado, en parte, ingiríendolas de forma sublingual o bucal (disolviendo la píldora debajo de tu lengua o entre tus encías y mejilla, respectivamente) en contraposición a oralmente (tragándola de forma normal). Ésto evita el metababolismo de primer paso en el hígado. Incluso con los métodos sublingual y bucal es común tragar alguna cantidad de la píldora, por lo que es posible asumir que existe algo de riesgo.
Por favor, comprende que el riesgo absoluto es de todas formas bajo (p.e., el \textit{paracetamol} es un orden de magnitud más preocupante con respecto al hígado que el estrógeno). Sin embargo \textbf{este riesgo se agrava con los riesgos del estrógeno relacionados con la nicotina.} Ver Pregunta \ref{11-2} también.

Más allá de esto, hay numerosos otros problemas de las píldoras que provienen de dos características principales: 1) Su corta vida media y biodisponibilidad pobre, y 2) su necesidad común de antiandrógenos. La primera característica hace que las píldoras sean mayormente no aptas para monoterapia (discutido luego) comparadas con las inyecciones. La segunda viene acompañada por diversos efectos secundarios negativos, dependiendo de los antiandrógenos involucrados (ver sección \ref{AA} “ANTIANDRÓGENOS”).  Juntas, estas características suman adicionales grados de variabilidad que generan regímenes pobres, y sus efectos secundarios (como poca energía/libido y resultados más lentos) más comunes que con otros tipos de administración. Las píldoras son además más dificiles de acumular, y, en algunos mercados, son más caras que los viales. Ten en cuenta que además, importar píldoras de distribuidores extranjeros en grandes volúmenes puede hacer sonar las alarmas aduaneras, lo que puede generar confiscaciones, pérdida monetaria y/o problemas legales, dependiendo de las leyes del país en que residas. \textbf{Si alguien pregunta, no sabes quien pidió esas píldoras.}

\textbf{Si por alguna razón consumes el estrógeno en píldoras, por favor toma 4-8 mg. sublingualmente de forma espaciada durante el día.} Por debajo de 4 mg. no es casi nunca una dosis suficiente.

\subsection{¿Qué tienen de malo los parches?}

\begin{itemize}
  \item Relativamente caros (típicamente incluso más que las píldoras);
  \item Más dificil de conseguir mediante DIY (a través de mercado gris);
  \item Generalmente necesitan de un antiandrógeno (ver sección \ref{AA} “ANTIANDRÓGENOS”);
  \item Pueden generar irritación de la piel;
  \item Requieren aplicación 24/7;
  \item Son propensos a caerse;
  \item No son siempre consistentes en su absorción (por ejemplo con temperaturas elevadas);
  \item Son difíciles de acumular (difíciles de adquirir de forma mayorista);
  \item A menudo es dificil subir de niveles de menopausia incluso con múltiples adheridos simultáneamente.
\end{itemize}

\subsection{¿Qué tiene de malo el gel?}

\begin{itemize}
  \item Difícil de dosificar de forma precisa, lo que lleva a niveles inconsistentes;
  \item Requiere una aplicación regular debido a su relativamente corta vida media;
  \item Puede ser un poco un engorroso (viscoso);
  \item Riesgo de exposición de segunda mano debido al contacto con otras personas
  \item Generalmente necesita un antiandrógeno (ver sección \ref{AA} “ANTIANDRÓGENOS”).
\end{itemize}

Cabe señalar, sin embargo, que el gel requiere mínimos recursos para la producción, lo cual puede ser una bendición en algunas circunstancias.

\subsection{¿Qué hay de los pellets?}

\begin{itemize}
  \item Generalmente mucho más caro que cualquier otra opción;
  \item Pocos proveedores los ofrecen;
  \item Los periodos de ajuste de dosis están muy diseminads temporalmente;
  \item Los pellets fallidos pueden generar niveles insuficientes;
  \item Pellets aplastados/rotos pueden causar niveles inesperadamente altos;
  \item Generalmente no es posible conseguirlos mediante DIY (mercado gris).
\end{itemize}

Particularmente el ultimo punto significa que sólo se pueden comprar de unos pocos proveedores probablemente caros. Es posible que esta sea la primera vez que escuchas de los pellets. ¿Ves el problema?

\subsection{¿Qué pasa con los sprays/aerosoles?}

Son bastante experimentales por lo que no hay mucho para decir de ellos, pero comparten los pros y contras del gel. Esta sección está principalmente para que puedas saber de su existencia.

\subsection{¿Es tan significante la diferencia?}

\textbf{Sí.} Hasta el punto que escribí todo este documento para poder enlazarlo y así tener que repetirme menos. Un régimen de inyecciones bien dosificado es la mejor forma de estrógeno que tenemos para alcanzar niveles de monoterapia. 

\subsection{¿Es precisa ésta tabla?}\label{1-12}

\begin{figure}[H]
  \centering
  \includegraphics[width=1\linewidth]{STUPID_CHART_evil_bad_bad_destroy_evil_bad.png}
  \caption{Esta tabla es basura}
  \label{fig:scebbdeb}
\end{figure}

\textbf{No.} Aunque esta guía no está interesada en responder individualmente a cada pedazo de desinformación de redes sociales, la prevalencia de esta tabla tanto online como en recursos provistos por doctores en diversos idiomas unida al gran volumen de daño que ha causado denota la necesidad de hacer una excepción especial más alla de una mera referencia de pasada en la Pregunta \ref{11-6}.

\textbf{Esta tabla es categóricamente falsa.} Casi todos los aspectos son engañosos en algún sentido, con la única excepcion de que es enteramente cierto que el estrógeno no genera cambios en la voz. Los rangos de tiempo en “expected onset” (momento previsto) son engañosos, y la idea de un límite de tiempo para el “maximum effect” (efecto máximo) es tan engañoso que bordea lo criminal. Una multitud de cambios no estan listados (como los efectos mentales) o son erróneos (la tabla se contradice a sí misma sobre la disfunción erectil). No se hace tampoco ninguna consideración con respecto a la calidad del régimen.

\textbf{Los aspectos clave de la tabla requieren un matiz que será discutido a lo largo de esta guía, por lo que la conclusión general es que los cambios suceden a diferentes ratios.} Esta tabla genera bajas espectativas, provee una imagen inexacta de lo que hace la TRH, y predispone a las personas trans al fallo limitando su entendimiento de hormonas. Por favor, haz caso omiso a la tabla entera. Continuemos ahora despues de este pequeño desvío.



\section{¿POR QUÉ INYECCIONES?}

\subsection{¿Qué hace a las inyecciones tan buenas?}

\textbf{Consistencia.} La consistencia es clave cuando se trata de la TRH. Tener un nivel consistente de hormonas significa estabilidad, y la estabilidad es buena. Incluso el "peor" tipo de inyección (sigue leyendo) puede proporcionar un ciclo hormonal más consistente que otras vías de administración, lo que proporciona muchos beneficios.

\subsection{¿Se necesitan antiandrógenos con las inyecciones?}

Generalmente no. Contar con un ciclo de inyecciones correctamente dosificado y espaciado que pueda proporcionar niveles de estrógeno lo suficientemente altos como para suprimir la producción de testosterona elimina la necesidad de un antiandrógeno en la mayoría de los casos. Esto se conoce como \textit{"monoterapia"}.

\subsection{¿Cómo funciona la monoterapia?}\label{2-3}

En términos simples, al cerebro le da igual qué hormona tiene, siempre y cuando tenga suficiente. Si hay consistentemente suficientes hormonas en tu cuerpo, deja de producir más. La parte "consistente" es lo que las inyecciones son capaces de hacer que a otras vías de administración les cuesta lograr. Intentar hacer una monoterapia suficiente con píldoras, por ejemplo, es probablemente imposible en la mayoría de las situaciones. En términos más específicos con respecto a \href{https://en.wikipedia.org/wiki/Hypothalamic-pituitary-gonadal_axis}{el eje HPG}, la \textit{hormona luteinizante} (LH) y la \textit{hormona folículo estimulante} (FSH) son suprimidas por niveles elevados de \textit{estradiol} (E2) en suero, lo que inhibe la producción de testosterona en los testículos.

\subsection{¿Por qué son más seguras las inyecciones?}

Al eliminar la necesidad de antiandrógenos (ver sección \ref{AA} “ANTIANDROGENS”), las inyecciones eliminan los riesgos a largo plazo asociados con los antiandrógenos. El estrógeno bioidéntico, al evitar el hígado (ver Pregunta \ref{11-1}), es lo más cercano que podemos estar a la producción natural de estrógeno, lo que elimina riesgos adicionales.

\subsection{¿Pero no hay riesgos con el acto físico de inyectarse?}

Sí, pero con una mínima práctica (ver Sección \ref{ts} “TÉCNICA Y MATERIALES”) en el peor de los casos puedes experimentar un pequeño moratón. Es similar a andar en bicicleta: Una vez que sabes cómo hacerlo, tendrías que esforzarte MUCHO para hacerlo significativamente mal.

\subsection{¿Por qué son más fáciles las inyecciones?}

Una vez que estás acostumbrada, son más fáciles. Las inyecciones no requieren una administración frecuente (por ejemplo, una inyección semanal contra múltiples píldoras diarias), no hay un gran riesgo de usar una dosis imprecisa (comparado con el gel), no se pueden despegar en medio de un ciclo (como los parches) y no requieren viajes potencialmente significativos a un proveedor (comparado con los pellets).

\subsection{¿Por qué son más económicas las inyecciones?}

En términos sencillos, se necesita mucho menos estrógeno. Un vial de 5 ml. que es capaz de proporcionar casi un año de estrógeno tiene solo 200 mg. de estrógeno en ese vial, mientras que un suministro mínimo equivalente de píldoras por ejemplo (4 mg. x 365 días = 1460 mg.) es sustancialmente mayor. Esto no es una comparación rigurosa, pero es una demostración útil de las proporciones. Otra comparación divertida es que puedes meter 1-2 años de viales de estrógeno dentro de un típico frasco de suministro de tres meses de píldoras.

\subsection{Pero no tengo seguro / mi seguro no lo cubre / las píldoras son más baratas que las inyecciones en mi caso / las inyecciones no están disponibles en mi país / mi médico no me receta inyecciones}

Por favor, mira la sección \ref{sv} “OBTENIENDO VIALES”. Te sorprenderás, y muy probablemente, te radicalizarás.

\subsection{¿Es buena idea cambiar a inyecciones incluso después de años con TRH?}

\textbf{Sí.} Nada está garantizado, pero muchas personas experimentan diferencias sustanciales y notables después de cambiar a inyecciones incluso después de años en la TRH. Estas van desde un aumento en el desarrollo mamario, mejor salud mental, alivio de los efectos secundarios de los antiandrógenos u otras formas de estrógeno, sentirse mejor en general, etc. Cambiar vale la pena.

\textbf{Debería mencionar que el tiempo en la TRH antes de las inyecciones no es tiempo "perdido", ni hay una ventana limitada en la que el estrógeno es más efectivo para la feminización.} Es la consistencia de las inyecciones lo que las hace tan buenas, pero el destino será en gran medida el mismo de cualquier forma. Ver Preguntas \ref{11-5} y \ref{11-6}.

\subsection{Pero las inyecciones dan miedo}

Sí, lo dan al principio. A nadie le gustan las agujas porque el cuerpo naturalmente no quiere pincharse a sí mismo, pero con la técnica y los suministros adecuados, no te dolerá apenas. Hay innumerables casos de personas con fobias debilitantes a las agujas que ahora les parece la experiencia de inyectarse aburrida. El miedo es normal y común, pero es completamente superable y vale la pena superarlo. "Oh, no era tan malo como pensaba" es una frase muy común. Como dice el mantra: Hazlo asustada. Estarás bien.

\subsection{¿Las inyecciones son como una extracción de sangre o una vacuna?}

No. Las extracciones de sangre típicamente usan agujas mucho más grandes y van a un lugar más sensible además de que también te sacan la sangre, lo cual suele ser desagradable. Las vacunas contienen ingredientes que causan reacciones inmunes dolorosas porque son vacunas. Las inyecciones de TRH ponen una pequeña cantidad de hormonas en ti. Confío en que ves la diferencia. El acto de inyectarte a ti misma también puede ser más fácil que el que otra persona te inyecte, dependiendo de tu inclinación.

\subsection{¿Hay alguna herramienta de accesibilidad para las inyecciones?}

Sí. Los auto-inyectores existen y pueden ser bastante útiles si tienes problemas de control motor fino, por ejemplo. Por favor, mira la Pregunta \ref{5-21}, o simplemente sigue leyendo.

\subsection{Pero \textit{soy} especial y no puedo inyectarme a mí misma porque tengo huesos de cristal y piel de papel y\textemdash{}}

Entiendo el miedo, pero si realmente no deseas hacer inyecciones bajo ninguna circunstancia y no tienes algún tipo de contraindicación legítima como la hemofilia, entonces no lo hagas. Puedes simplemente decirlo. Está bien. Cuando cambies de opinión, esta guía seguirá aquí. Y si no lo haces, que así sea.



\section{TIPOS Y DOSIS}\label{td}

\subsection*{Vocabulario Clave}
\addcontentsline{toc}{subsection}{\textemdash{} Vocabulario Clave}

\subsection{¿Cuáles son los diferentes tipos de estrógeno inyectable?}

Los cuatro tipos principales usados para la TRH son el \textit{Valerato de Estradiol} (EV, del inglés \textit{"Estradiol Valerate"}), el \textit{Cipionato de Estradiol} (EC, del inglés \textit{"Estradiol Cypionate"}), el \textit{Enantato de Estradiol} (EEn, del inglés \textit{"Estradiol Enanthate"}) y el \textit{Undecanoato de Estradiol} (EUn, del inglés \textit{"Estradiol Undecylate"}). Cada uno de estos es un “éster” de \textit{estradiol} (E2) y será convertido a \textit{estradiol} en tu cuerpo.

Por favor, ten en cuenta que en algunas regiones se venden píldoras confusamente con el nombre de \textit{Valerato de Estradiol}, pero esta sección solo se refiere a la forma inyectable.

\subsection{¿Cuáles son las diferencias entre cada tipo de estrógeno inyectable?}

La única diferencia relevante entre ésteres es que cada uno tiene una vida media diferente y una curva hormonal resultante diferente, lo que a su vez afecta la dosis y la frecuencia.

\subsection{¿Feminiza mejor un tipo de estrógeno inyectable que otro?}

\textbf{No.} Las diferencias afectan a la dosis y a la frecuencia, lo que constituye una diferencia cualitativa en la experiencia, lo cual puede hacer que un éster sea preferible a otro, pero los cuatro tipos funcionan aceptablemente bien y conservan los mismos beneficios del resto de inyecciones.

\subsection{¿Qué tipo de estrógeno inyectable debería elegir si tengo la opción?}

Si tienes la opción, el \textit{Enantato de Estradiol} se prefiere para la mayoría de las personas debido a los niveles excepcionalmente estables que proporciona, con la salvedad de que en la mayoría de los países esta opción solo existe si estás usando el mercado gris (DIY, ver Sección \ref{sv} “OBTENIENDO VIALES”). Si estás yendo a través de un doctor, es posible que tengas la opción del \textit{Cipionato de Estradiol}, pero usualmente en bajas concentraciones lo que puede hacer que los beneficios sean nulos dependiendo de tu tolerancia a las inyecciones de alto volumen. El éster inyectable más comúnmente recetado, el \textit{Valerato de Estradiol}, sigue siendo completamente capaz de producir buenos resultados, pero tiene algunas pequeñas molestias que lo hacen no preferible cuando hay otra opción (es decir, cuando usas DIY). Sigue leyendo.

\subsection{¿Qué significa "concentración"?}

Los viales de estrógeno están hechos de estrógeno contenido en una solución de aceite. La \textit{concentración} de un vial es la cantidad de estrógeno contenida en esa solución. Esto se da como una proporción de masa a volumen para el vial. En otras palabras: Por cada mililitro de aceite (medida de volumen), hay esa cantidad de miligramos de estrógeno (medida de masa). \textbf{Por favor, entiende que la concentración por sí sola no es una dosis.} A menudo verás las concentraciones listadas por el volumen total del vial (p.ej., 200 mg/5 ml) pero siempre es preferible simplificar esta fracción (40 mg/ml en este caso). \textbf{Las concentraciones típicas son 5 mg/ml, 10 mg/ml, 20 mg/ml, 40 mg/ml y ocasionalmente 50 mg/ml.}

\subsection{¿Qué significan "dosis" y "frecuencia"?}

\textit{dosis} y \textit{frecuencia} son los dos factores que determinan tu ciclo hormonal. \textit{dosis} se refiere a cuánto estrógeno pones en ti (medido en mg.), y \textit{frecuencia} se refiere a cómo de seguido pones estrógeno en ti (medido en días o semanas). A menudo escucharás la palabra “régimen” también, que se refiere a todo lo relacionado con la TRH que estás tomando y en qué frecuencias.

\subsection{¿Cómo sé cuál es mi dosis?}

Tu dosis es la concentración de tu vial multiplicada por el volumen que estás inyectando. \[Concentración (mg/ml) x volumen (ml) = dosis (mg)\] \textbf{Por favor, entiende que el volumen por sí solo no es una dosis.} Una analogía sería con la cocina: No puedes simplemente decir “hornea por 45 minutos” porque tienes que saber a qué temperatura poner el horno.

\subsection{¿Cuáles serían ejemplos típicos de cálculos de dosificación?}

Las matemáticas son simples, ¡lo prometo! A continuación hay una pequeña tabla de referencia comparando concentraciones y volumen para una gama de dosis comunes. Limítate a solo dos decimales. No estarás usando jeringas que tengan la precisión para un número como 0.153 ml. por ejemplo. Eso está dentro del error de redondeo y no es una diferencia relevante para nuestra escala. También puedes usar \href{https://hrtcafe.net/Calc/}{esta calculadora} como referencia también.

\begin{table}
\centering
\caption{Dosis de ejemplo para concentraciones comunes por volumen}
\label{tab:concentrations}
\begin{tabular}{@{}llllll@{}}
  \toprule
  \multicolumn{1}{c}{} & \multicolumn{5}{c}{Concentraciones (mg/ml)} \\
  \cmidrule(rl){2-6}
            & 5    & 10  & 20 & 40 & 50   \\
            \cmidrule(rl){2-6}
dosis (mg) & \multicolumn{5}{c}{Volumen (mL)}  \\
    \cmidrule(r){1-1} \cmidrule(lr){2-6}
4        & 0.8  & 0.4 & 0.20  & 0.10 & 0.08\\
5        & 1.0    & 0.5 & 0.25 & 0.13 &  0.10\\
6        & 1.2  & 0.6   & 0.30  & 0.15 & 0.12   \\
7        & 1.4  & 0.7 & 0.35  & 0.18 & 0.14\\
8        & 1.6  & 0.8   & 0.40  & 0.20 & 0.16  \\
9        & 1.8  & 0.9 & 0.45  & 0.23 & 0.18\\
10       & 2.0    & 1.0   & 0.50  & 0.25 & 0.20 \\
  \bottomrule
\end{tabular}
\end{table}

\textbf{Cómo leer esta tabla:} Toma tu dosis deseada a la izquierda y encuentra el volumen correspondiente a la derecha para tu concentración dada en la columna superior. Notarás que los requerimientos de volumen para que los viales de 5 mg/ml tengan dosis razonables no son buenos. Eso es porque los viales de 5 mg/ml no son buenos.

\subsection{¿Cómo convierto dosis entre ésteres?}

\textbf{No puedes.} Ya que se comportan de forma diferente, no existe una “conversión” entre dosis en ese sentido. Si cambias de un éster a otro, simplemente debes hacer una dosis típica para el nuevo éster y trabajar desde ahí. Puedes hacer comparaciones entre ellos, pero no hay un método para convertir de uno a otro.

\subsection{¿Cómo puedo comparar diferentes curvas y dosis entre ésteres?}

Si te gustaría meterte en terreno "nerd", tengo en muy alta consideración a \href{https://estrannai.se}{estrannai.se}. Ten en cuenta que esto no es obligatorio de usar pero es una buena herramienta para realizar comparaciones aproximadas. \href{https://estrannai.se/\#i0__cu,7,7,1-cu,5,7,3-cu,5,7,2}{Aquí hay un ejemplo de comparación entre dosis semanales típicas} que ahora veremos individualmente.

\textbf{Hay que tener en cuenta que las dosis que enumero a continuación deberían ser suficientes en el extremo inferior del rango en la mayoría de los casos.} Comienza con el número más bajo y sube si lo necesitas. Más no es inherentemente mejor, pero hablaremos de eso en profundidad más adelante. Estos rangos de dosis son poco probables que cambien independientemente de dónde hayas adquirido tu vial.

\subsection*{Conoce a tus ésteres}
\addcontentsline{toc}{subsection}{\textemdash{} Conoce a tus ésteres}

\subsection{¿Cómo dosifico el \textit{Valerato de Estradiol}?}

O bien dos veces por semana a una dosis más baja o una vez por semana a una dosis más alta es necesario para buenos niveles con \textit{Valerato de Estradiol}. Es una cuestión de comodidad y tolerancia. La regla general típica es aproximadamente 1 mg. por cada día en un ciclo con frecuencias generalmente entre 3-7 días. \textbf{Una dosis semanal entre 6-8 mg. es mi recomendación típica}, pero 4-5 mg. por 5 días también es muy común. \textbf{La frecuencia nunca debe ser menos frecuente que semanal (es decir, no más de siete días entre inyecciones).} Semanal ya está empujando cuánto puede durar el éster. Cualquier cosa más allá de eso es altamente desaconsejada para evitar efectos secundarios relacionados con la fluctuación (ver Pregunta \ref{7-3}).

Por favor, ten en cuenta que en algunas regiones se venden píldoras confusamente con el nombre de \textit{Valerato de Estradiol}, pero esta sección solo se refiere a la forma inyectable.

\subsection{¿Cómo es la curva hormonal del \textit{Valerato de Estradiol}?}

El \textit{Valerato de Estradiol} es el éster más quisquilloso de los cuatro. Rápidamente alcanza un pico muy alto unos días después de la inyección y tan rápidamente cae de nuevo. Esta relativa inestabilidad puede ser desagradable dependiendo de tus sensibilidades personales, pero con ajustes a la frecuencia y dosis esto se puede mitigar hasta cierto punto.

\begin{figure}[H]
  \centering
  \includegraphics[width=1\linewidth]{ev.png}
  \caption{Serum Estradiol (pg / ml) of Estradiol Valerate vs Time (days)}
  \label{fig:ev}
\end{figure}

\subsection{¿Cómo dosifico el \textit{Cipionato de Estradiol}?}

El \textit{Cipionato de Estradiol} puede acomodar una dosis semanal sin problema. \textbf{Una dosis semanal entre 5-7mg es típica.} Extender la duración más allá de lo semanal (p.ej., cada 10 días) no se recomienda porque es un uso menos eficiente del estrógeno en comparación con las dósis semanales ya que requiere dosis cada vez más altas para alcanzar niveles aceptables. Cualquier extensión más allá de lo semanal es mucho más propensa a efectos secundarios debido a la fluctuación (ver Pregunta \ref{7-3}).

\subsection{¿Cómo es la curva hormonal del \textit{Cipionato de Estradiol}?}

El \textit{Cipionato de Estradiol} es un poco más indulgente que el \textit{Valerato de Estradiol}. La curva no progresa tan rápidamente y tiene una variación mucho menor entre el pico y el valle, pero aún así hay una subida y bajada notable durante un ciclo típico semanal.

\begin{figure}[H]
  \centering
  \includegraphics[width=1\linewidth]{ec.png}
  \caption{Estradiol Suero (pg/ml) del Cipionato de Estradiol vs Tiempo (días)}
  \label{fig:ec}
\end{figure}

\subsection{¿Cómo dosifico el \textit{Enantato de Estradiol}?}

El \textit{Enantato de Estradiol} puede facilmente acomodar una dosis semanal sin problema y posiblemente se puede extender hasta 10 días si así lo deseas. Más allá de eso es técnicamente posible pero no recomendado ya que los niveles se volverán cada vez más inestables. \textbf{Una dosis semanal de 4-6mg es típica}, con 5-7mg recomendados para hasta 10 días. Semanal sigue siendo recomendado independientemente para mayor consistencia y facilidad de planificación, ya que cualquier extensión hasta 10 días no ofrece mucho beneficio en mi opinión.

\subsection{¿Cómo es la curva hormonal del \textit{Enantato de Estradiol}?}

El \textit{Enantato de Estradiol} es el standard de oro para el estrógeno inyectable. Tiene una curva que es extremadamente plana (es decir, tiene poca fluctuación) durante la duración típica de una semana. Esto permite niveles muy consistentes sin efectos secundarios negativos relacionados con la fluctuación (ver Pregunta \ref{7-3}).

\begin{figure}[H]
  \centering
  \includegraphics[width=1\linewidth]{een.png}
  \caption{Estradiol Suero (pg/ml) del Enantato de Estradiol vs Tiempo (días)}
  \label{fig:een}
\end{figure}

\subsection{¿Cómo dosifico el \textit{Undecanoato de Estradiol}?}

El \textit{Undecanoato de Estradiol} es capaz de extenderse mucho más allá de lo semanal, llegando al rango mensual o trimestral. La dosificación recomendada para esto, sin embargo, no es estandarizada ni conocida. Los factores que afectan cómo se absorbe el estrógeno de una inyección (\textit{“farmacocinéticos”}) que son insignificantes para otros ésteres, son significativos para el \textit{Undecanoato de Estradiol}. Por tanto, esto sigue siendo un territorio altamente experimental que está más allá del alcance de esta guía. Considera consultar un almanaque de brujas para el calendario lunar para inyectarte una vez cada luna llena.

\subsection{¿Cómo es la curva hormonal del \textit{Undecanoato de Estradiol}?}

Realmente no lo sabemos. Los datos son demasiado escasos para pintar una imagen precisa de su curva en su totalidad, y las variables son abundantes. Es algo que puedes investigar y experimentar si estás interesada, pero es un terreno nuevo y necesitas entender los riesgos que implica ser un conejillo de indias humano, así que no lo recomiendo a menos que sepas lo que estás haciendo.

\begin{figure}[H]
  \centering
  \includegraphics[width=1\linewidth]{moon.png}
  \caption{La Luna}
  \label{fig:moon}
\end{figure}



\section{ANÁLISIS DE SANGRE Y NIVELES}

\subsection*{Obteniendo Resultados}
\addcontentsline{toc}{subsection}{\textemdash{} Obteniendo Resultados}

\subsection{¿Cada cuánto debo testear mis niveles?}

Mientras estás ajustando tu dosis por primera vez, querrás testear de forma relativamente frecuente. Después de cualquier ajuste a tu régimen, deberías darle a tus niveles 1-2 meses para estabilizarse, y luego testear una vez que hayan alcanzado su nuevo normal.

\subsection{¿Tengo que testear mis niveles antes de comenzar la TRH?}

Posiblemente no, porque la testosterona estará demasiado alta y el estrógeno demasiado bajo, lo cual no son datos particularmente útiles. Sin embargo, un análisis de sangre general rutinario (es decir, un panel de lípidos y demás) es recomendable para tu salud de todas formas. La excepción es si crees que puedes tener una condición intersexual que pueda afectar tu régimen de TRH ya que a veces esto puede ser visible en el análisis de sangre preliminar.

\subsection{¿Tengo que testear mis niveles si no he cambiado mi dosis en un tiempo?}

Probablemente no, porque si no has cambiado nada entonces nada debería haber cambiado. Puede ser bueno para tu tranquilidad si has cambiado aspectos de tu régimen/proveedor, y los doctores/seguros a menudo lo requieren, pero no se debería esperar una desviación importante. Una salvedad es que si estás experimentando con \textit{Undecanoato de Estradiol}; en ese caso, casi con seguridad deberías testear trimestralmente como mínimo de todas formas.

\subsection{No tengo seguro o doctor. ¿Dónde puedo obtener un análisis de sangre?}

Busca opciones de análisis de sangre privadas en tu región. En muchas ubicaciones puedes legalmente obtener análisis de sangre privados, pero podrían no ser baratos. Puede haber opciones en línea que te permitan obtener esos análisis con algún descuento pero depende mucho de tu región.

\subsection{No puedo obtener/permitirme un análisis de sangre. ¿Todavía puedo hacer TRH?}

Aunque tener la información es obviamente preferible a no tenerla, la TRH es extremadamente segura y a dosis típicas no debería representar ningún problema. Simplemente tendrás que confiar más en cómo te sientes y lo que observas.

\subsection{¿Qué niveles debería testear?}

Al menos el \textit{Estradiol} (E2) y la \textit{Testosterona Total} (T) porque son los niveles principales de las que debes preocuparte. La \textit{Globulina Fijadora de Hormonas Sexuales} (SHBG, del inglés \textit{"Sex Hormone Binding Globulin"}), la \textit{Dihidrotestosterona} (DHT), la \textit{Estrona} (E1), y la \textit{Prolactina} (PRL) también pueden ser útiles para testear si estás experimentando problemas porque pueden servir para tratar de encontrar soluciones. La \textit{Hormona Foliculoestimulante} (FSH, del inglés \textit{"Follicle-Stimulating Hormone"}) y la \textit{Hormona Luteinizante} (LH, del inglés \textit{"Luteinizing Hormone"}) pueden decirte si tu eje \textit{Hipotálamo-Hipofisario-Gonadal} (HPG, del inglés \textit{"Hypothalamic-Pituitary-Gonadal"}) está inactivo, lo cual es la base de la monoterapia (ver Pregunta \ref{2-3}). Pero de nuevo: El \textbf{\textit{Estradiol} y la \textit{Testosterona Total} son las principales preocupaciones.}

\subsection{¿Cuándo debería hacerme un análisis de sangre durante mi ciclo hormonal?}

Al final de tu ciclo (su momento \textit{valle}, o \textit{“trough”} en inglés). Quieres estar lo más cerca posible del fondo porque esta es la información más útil. Posiblemente sea la única información útil ya que tener niveles mínimos consistentes es la preocupación principal. Por ejemplo, si normalmente te inyectas el jueves por la tarde, hazte tus análisis el jueves siguiente por la mañana o temprano por la tarde antes de tu próxima inyección.

\subsection{Mi doctor dijo que debía hacer análisis de sangre en el punto medio/en el pico, ¿debería hacerlo?}

\textbf{No.} Medir el nivel máximo de estrógeno no proporciona información útil y es solo una medida de qué éster estás usando. Siendo benévolas, podemos pensar que es incompetencia debida a estándares de cuidado conservadores y desactualizados. Si no lo somos, podríamos pensar que es malicia para asegurar niveles insuficientes de estrógeno que resultarán en mala salud, resultados lentos o cualquier otro tipo de resultados negativos. \textbf{Recomiendo medir en el valle de todas formas.}

\subsection*{Interpretando Resultados}
\addcontentsline{toc}{subsection}{\textemdash{} Interpretando Resultados}

\subsection{¿Qué niveles de estrógeno quiero?}

Esta es probablemente la pregunta más controvertida en el tema de la transición. La respuesta corta es que quieres suficiente para sentirte bien y para suprimir la testosterona si lo necesitas, pero más allá de eso, niveles más altos son innecesariamente derrochadores en el mejor de los casos y pueden ser contraproducentes en el peor de los casos. Sin embargo, esto es un rango amplio, y con tantas variables siempre hay una desviación personal. En otras palabras: Quieres suficiente estrógeno para sentirte bien, y eso es todo.

\subsection{¿Los niveles más altos de estrógeno feminizan mejor o más rápido?}

\textbf{No.} Niveles de estrógeno más altos de lo necesario podrían ser preferidos por alguien por su experiencia subjetiva, pero no confieren beneficios de feminización. De hecho, unos niveles demasiado altos pueden sentirse mal al causar inestabilidad del estado de ánimo u otros efectos secundarios indeseables. \textbf{Minimizar los niveles de testosterona a un nivel basal es mucho más importante para la feminización que maximizar los niveles de estrógeno.}

\subsection{De acuerdo, ¿pero qué número quiero ver en los resultados de estrógeno de mi analítica?}

Con el entendimiento de que el número exacto no importa, que el número siempre será un poco más alto que lo que hay en tu cuerpo incluso en un día de valle debido a la latencia, y que el número estará en una nube de posibilidades basada en cualquier cantidad de factores, \textbf{recomiendo un valle de aproximadamente 200 pg/ml (730 pmol/L) como mínimo.} Esta es una recomendación ligeramente conservadora para proporcionar un amplio margen de maniobra ya que la supresión del eje HPG ocurre por debajo de esto. Alrededor de ese valor tiende a funcionar bien para la mayoría, aunque algunas prefieren más alto o más bajo. No creo que este sea un número al que se le deba prestar demasiada atención porque es inherentemente variable y si te sientes bien eso es lo que más importa, \textbf{pero más allá de 300 pg/ml (1100 pmol/L) en el valle es casi con seguridad un valor más alto de lo que necesita ser o debería ser.} Hay excepciones a esto, pero probablemente no seas la excepción. Aún así, haz lo que te parezca correcto. Además, mira la Pregunta \ref{11-1}.

\subsection{¿Qué niveles de testosterona quiero?}

La supresión de la testosterona es el requisito clave para una feminización adecuada, por lo que menos de 50 ng/dL (1.7 nmol/L) es generalmente suficiente. \textbf{Notablemente, la testosterona cercana a cero no es deseada.} Ver Sección \ref{T} “TESTOSTERONA”.

\subsection{Tengo testosterona alta/baja de forma natural. ¿Necesito ajustar mi dosificación?}

Probablemente no. El rango de testosterona que se encuentra típicamente antes de empezar la TRH es casi siempre más alto de lo que se desea para la feminización y aún así será suprimido en cualquier caso (ver Pregunta \ref{2-3}). La excepción sería si tienes alguna variedad de condiciones intersexuales que puedan causar la necesidad de un ajuste más fino que las recomendaciones listadas en esta guía, lo cual está más allá del alcance de lo que esta guía puede proporcionarte. Puede que no necesites ajustar, pero tal vez te sientas mejor si lo haces. En última instancia, haz lo que te parezca correcto. Mira también la Pregunta \ref{9-2}.

\subsection{He tenido cirugía de reasignación. ¿Mis niveles de estrógeno necesitan ser diferentes?}

Dado que la supresión de la testosterona ya no es una preocupación para ti, es posible que aún puedas sentirte genial con niveles de estrógeno más bajos que los que tienes actualmente, \textbf{pero aún necesitas estrógeno.} Debido a que ya no produces tus propias hormonas, es crucial que mantengas niveles suficientes de hormonas para tu salud. Tener pocas o ningunas hormonas conducirá a síntomas de menopausia, que es la misma razón por la que las mujeres cis mayores podrían tomar TRH una vez que llegan a la menopausia. Ajusta según lo veas necesario.

Para claridad adicional: \textbf{Mantener un mínimo de aproximadamente 100 pg/ml (350 pmol/L) es esencial para evitar preocupaciones de densidad mineral ósea.} Si la mayor parte de tu feminización ya está completa, entonces en muchos aspectos tu perfil hormonal es comparable al de una mujer cis menopáusica, por lo que se pueden aprender lecciones de ellas (ver Pregunta \ref{11-29}). En algunos casos de fatiga o baja energía, suplementar con dosis bajas de testosterona puede ser beneficioso (ver Pregunta \ref{9-2}).

\subsection{¿Hay algo que pueda causar que una prueba de sangre sea inexacta?}

Dependiendo de cómo se mida la sangre (el \textit{“ensayo”}), los suplementos de biotina pueden causar que los niveles de \textit{estradiol} (E2) sean inesperadamente altos (entre otros, pero el \textit{estradiol} es nuestra preocupación). No siempre es posible saber el tipo de ensayo que se usará, por lo que se recomienda pausar cualquier suplemento de biotina unos días antes de las pruebas. También es posible que haya habido un error con el equipo o la muestra, aunque esto no es probable.

\subsection{¿Se muestran los diferentes ésteres de estrógeno o vías de administración de manera diferente en las pruebas de sangre?}\label{4-16}

No. No hay forma de saber qué tipo de éster de estrógeno o vía de administración alguien está usando basándose únicamente en el resultado de un análisis de sangre. Los diversos ésteres inyectables se convierten todos en \textit{estradiol} como queremos, y lo mismo ocurre con las píldoras, parches, geles, sprays o cualquier otra cosa. Al final del día, todo es estrógeno.



\section{TÉCNICA Y SUMINISTROS}\label{ts}

\subsection*{Sitios \& Seguridad}
\addcontentsline{toc}{subsection}{\textemdash{} Sitios \& Seguridad}

\subsection{¿Cómo hago una inyección de manera segura?}

Recomiendo los siguientes dos videos (en inglés):

\begin{enumerate}
  \item \href{https://www.youtube.com/watch?v=cBabaGC2Dok}{\textit{“¿Cómo realizar una inyección intramuscular (IM) por uno mismo?”}}
  \item \href{https://www.youtube.com/watch?v=YfNlAZLxLyw}{\textit{“Técnica de inyección IM sin dolor (para mí hasta ahora)”}}
\end{enumerate}

Entre estos dos videos, deberías estar completamente equipada para inyectarte de manera adecuada con un dolor mínimo. Te sugiero estudiarlos y revisarlos según necesites. \textbf{Una cosa clave a enfatizar es inyectar con la bisel hacia arriba para reducir el dolor.} En otras palabras: la aguja tiene un punto claramente definido, y quieres que eso sea lo que toque tu piel primero. Quieres que viaje en línea recta. Puedes pensar en cómo tu mano/muñeca rota si eso te ayuda a visualizar el movimiento, pero siendo realistas será memoria muscular intuitiva que aprenderás naturalmente.

\textbf{Recuerda, ¡inyectarte es una habilidad!} Mejorarás con el tiempo, y no te tomará mucho. ¡Tú puedes!

\subsection{¿Tengo que inyectarme exactamente así?}

No, la variación está bien. En definitiva, cuando la tarea es simplemente pincharte a ti misma, hay muchas formas de hacerlo. Encuentra la forma que mejor te funcione. Hacer un movimiento rápido generalmente funciona mejor, pero si tienes que ir lento, también funciona bien si es algo en lo que puedes mejorar de forma consistente.

\subsection{¿Cómo superar la ansiedad por la inyección?}

Sugiero transformar el proceso en un ritual. Al formar una rutina, el proceso se vuelve algo natural. Si puedes distraer tu mente escuchando música, teniendo una conversación, viendo un programa de televisión, o haciendo algo más que funcione para ti para dejar que tu memoria muscular tome el control, ¡eso es genial! Encuentra lo que te funcione a ti. Tener un amigo o ser querido que haga tus primeras pocas inyecciones también puede ayudar. Para la mayoría de las personas, la primera inyección es la más aterradora. Por lo general, la gente dice "oh, ¿eso fue todo?" porque nunca es tan malo como esperaban.

\subsection{¿Es importante dónde inyecto en mi cuerpo?}

Sí y no. Mantenerse dentro de áreas seguras es importante, pero fuera de eso, dónde inyectas depende principalmente de tu movilidad, del volumen de fluido que estás inyectando, la combinación de aguja/jeringa que estás usando y tu propia comodidad. De cualquier manera, \textbf{asegúrate de rotar los sitios de inyección.} Alterna los lados de tu cuerpo con cada inyección. Por ejemplo, si inyectas en tu pierna derecha una semana, usa tu pierna izquierda la siguiente. Esto es para minimizar cualquier riesgo a largo plazo de cicatrices.

\subsection{¿Qué sitios de inyección son seguros?}\label{5-5}

Las opiniones varían entre las autoridades médicas, pero tu composición corporal también puede jugar un papel. Recomiendo inyectar en un lateral de la pierna como se muestra en los vídeos porque es factible para la mayoría de las personas y es algo que permite consistencia, lo que significa inyecciones consistentemente indoloras una vez que hayas practicado tu técnica, pero otras personas prefieren su glúteo o su estómago. \href{https://vertisis.com/articles/how-to-self-administer-a-subcutaneous-injection}{Este video y sitio web} muestra otros sitios de inyección que pueden ser aceptables dependiendo de los suministros que uses. Descubre qué funciona mejor para ti.

Una cosa a tener en cuenta: Si eliges inyectar en tu estómago, es crucial que no inyectes demasiado profundo ya que eso será más doloroso.

\subsection{¿Qué significan "intramuscular" (IM) y "subcutánea" (SubQ/SC)?}

A menudo escucharás estos términos en el contexto de inyecciones. \textit{Intramuscular} significa inyectado en el músculo y \textit{subcutánea} significa inyectado en la capa grasa debajo de tu piel.

\subsection{¿Cuál es la diferencia entre inyecciones intramusculares (IM) y subcutáneas (SubQ/SC)?}

\textbf{En el contexto de la TRH, hay poca o ninguna diferencia entre las inyecciones intramusculares y subcutáneas.} Las inyecciones subcutáneas se absorben más lentamente que las intramusculares. Sin embargo, esta diferencia generalmente no es lo suficientemente significativa como para afectar a la dosificación. También debe tenerse en cuenta que una inyección raramente se deposita completamente en el músculo o completamente en la capa subcutánea, lo que difumina aún más cualquier diferencia cuando estamos hablando de múltiples inyecciones.

Como nota al margen, las fuentes farmacéuticas para los viales de estradiol típicamente dicen que son solo para inyecciones intramusculares porque eso es para lo que están técnicamente aprobados, pero no hay una diferencia inherente. El contenido del vial a menudo es idéntico de todos modos, y el contenido del vial es lo que importa en última instancia. Sigue leyendo.

\subsection{¿Debería realizar inyecciones intramusculares (IM) o subcutáneas (SubQ/SC)?}

\textbf{Ésta es la pregunta equivocada. Una inyección es una inyección.} Las inyecciones subcutáneas a menudo se recomiendan porque la gente cree que son menos dolorosas por el hecho de ser subcutáneas, pero no hay una diferencia fundamental en cómo se realiza una inyección. \textbf{Las ventajas a las que se refieren las personas no son inherentes a la ubicación del depósito de la inyección; son inherentes a los factores que afectan el dolor de la inyección.} La mejor pregunta sería "¿Cómo minimizo el dolor durante la inyección?", pero dos preguntas más primero.

\subsection{¿Importa mi ángulo de inyección y/o método preferido de inyección?}\label{5-9}

No. Reitero: La parte más importante de realizar una inyección es que atravieses tu piel con una aguja y deposites fluido en tu cuerpo. Si el fluido no se filtra (o al menos, no mucho) y no duele (o al menos, no mucho), entonces has hecho un trabajo fantástico. \textbf{No puedo enfatizar lo suficiente que la "división" entre intramuscular vs subcutánea es inexistente y que la pregunta no impacta significativamente la efectividad del estrógeno inyectable.} El \textit{Undecanoato de Estradiol} es el único caso donde la ubicación del depósito parece afectar significativamente a la absorción, pero incluso aquí no debería preocuparnos porque en este caso no entendemos completamente los detalles. La cuestión es: Por favor, preocúpate por las cosas que importan y no por las cosas que no importan.

La única consideración donde el ángulo importa sería si eliges hacer inyecciones en el estómago, ya que es fácil ir demasiado profundo si vas directamente hacia adentro, dependiendo de la longitud de tu aguja (discutida a continuación). Sin embargo el principio sigue siendo el mismo.

\subsection{¿Tengo que aspirar?}

No. La "aspiración" se refiere a tirar del émbolo hacia atrás después de perforar la piel antes de inyectar el fluido con la intención de asegurarse de que no se está inyectando en un vaso sanguíneo. Su necesidad es controvertida, pero para las inyecciones hormonales siguiendo procedimientos estándar hay pocos beneficios que superen sus contras. Los sitios de inyección standard tienen un bajo riesgo de tocar un vaso sanguíneo en primer lugar, disminuido aún más por las longitudes de aguja más cortas, así que esta práctica ya no es recomendada por la mayoría de las organizaciones médicas.

\subsection{¿Cómo minimizo el dolor durante la inyección?}

Además de practicar tu técnica y mejorar tu habilidad, el factor principal para la incomodidad de la inyección es la combinación de aguja y jeringa que estás usando. \textbf{Para minimizar la incomodidad, se debe usar la aguja de mayor calibre que el aceite portador de tu vial pueda tolerar junto con una jeringa y longitud de aguja apropiadamente dimensionadas.} Deberías preguntarte "¿Con qué calibre y longitud de aguja me debería inyectar?" Para responder a eso, hablemos de cómo funcionan las agujas.

\subsection*{Conociendo Tus Agujas}
\addcontentsline{toc}{subsection}{\textemdash{} Conociendo Tus Agujas}

\subsection{¿Qué es el "calibre de la aguja"?}

El \textit{calibre de la aguja} es una medida del grosor de la aguja. A mayor número, más fina es la aguja. Por ejemplo, una aguja 25G es más fina que una aguja 20G. Las agujas de mayor calibre tambien suelen tender a ser más cortas, porque las más largas tienen tendencia a doblarse, por lo que su largo tiene un máximo. Para sorpresa de nadie, las agujas más finas suelen doler menos. Es importante aclarar que el calibre de la/s aguja/s no va a afectar el TRH de ninguna manera, sólo va a afectar la facilidad y comodidad de la inyección en sí misma.

\subsection{¿Qué son las jeringas/agujas "luer lock" y "de insulina"?}\label{5-13}

Las \textit{jeringas luer lock} tienen jeringas y agujas separadas, por lo que se puede usar una aguja separada para extraer y otra para inyectar. (\textit{"luer lock"} es un sistema standard de bloqueo en el que la aguja se enrosca en la jeringa). Las \textit{jeringas de insulina} tienen una aguja fija, lo que significa que la misma aguja se usará para extraer e inyectar. Nota: No quieres agujas \textit{"luer slip"} (sistema de bloqueo donde la aguja se ajusta a presión) ya que son propensas a errores. Donde sea posible, las jeringas de insulina son preferidas por comodidad y por minimizar el espacio muerto (ver Pregunta \ref{5-26}).

\textbf{Advertencia de seguridad: Volver a colocar la tapa en las agujas generalmente no se recomienda por el riesgo de pincharte a ti misma, pero si lo haces (al igual que al cambiar una aguja de extracción), NUNCA apliques fuerza con tu mano hacia la aguja.}

Es posible que la tapa se rompa y puedas lastimarte si colocas la tapa incorrectamente. Se recomienda colocar la tapa cerca de la aguja en una superficie horizontal y presionar la aguja (o jeringa), aún con la tapa suelta, contra una pared o agarrar la tapa por los lados para que quede bien ajustada. No hay riesgo de transferencia de enfermedades al realizar una auto inyección, así que ten en cuenta esta advertencia según tu propio criterio, pero volver a colocar la tapa es una preocupación MUY seria al realizar inyecciones en otros. Para su deshecho, ver Pregunta \ref{5-27}.

\subsection{¿Qué calibre de aguja debo usar para extraer?}

Si estás usando jeringas \textit{luer lock} (enroscables), se recomienda usar un calibre más bajo (más grueso) que el que usas para inyectar, de forma que tome menos tiempo extraer del vial. Si es demasiado bajo puede causar \textit{coring} (ver Pregunta \ref{5-23}), así que se recomienda \textbf{extraer con al menos 21-23G}. Si tienes paciencia y volúmenes bajos que inyectar, entonces se recomiendan calibres más altos (menos gruesos) para la reducción del riesgo de \textit{coring} antes mencionado. Ten en cuenta que la aguja no se desafila significativamente en el tapón. Esta pregunta es irrelevante con las jeringas de insulina porque la aguja no es extraíble.

\subsection{¿Qué longitud de aguja debería usar para extraer?}

Si estás usando jeringas \textit{luer lock} (enroscables), la longitud de la aguja de extracción no importa mucho, fuera de la incomodidad de tener una aguja demasiado larga que sea difícil de manejar. En otras palabras, no es necesario ser exigente. Esta pregunta es irrelevante con las jeringas de insulina porque la aguja no es extraíble.

\subsection{¿Qué calibre de aguja debo usar para inyectar?}\label{5-16}

Esta es una pregunta difícil y altamente subjetiva, y tu respuesta dependerá de cuatro factores principales: 1) El aceite portador para lo que estás inyectando; 2) si el vial contiene un cosolvente; 3) tu paciencia para tener una aguja en tu cuerpo por más tiempo; y 4) tu disposición/capacidad para empujar más fuerte el émbolo de la jeringa. Es una cuestión de comodidad. Los aceites más gruesos significan más tiempo y más esfuerzo al usar un calibre más alto, pero también los calibres más altos pueden ser significativamente menos dolorosos al empezar. \textbf{Como punto de partida, 25G es el calibre mínimo de aguja que deberías usar para manejar la incomodidad.} La mayoría de los aceites portadores comunes generalmente pueden tolerar hasta 27G cómodamente, mientras que el aceite MCT en particular es notable por poder fácilmente tolerar hasta 30G (ver Pregunta \ref{6-16}).

\subsection{¿Con qué longitud de aguja debería inyectar?}\label{5-17}

\textbf{Recomiendo entre 0.5” a 1” (12.5mm a 25mm) dependiendo de tu calibre.} Por debajo de 0.5” (12.5mm) aumenta la probabilidad de fugas. Las agujas de 0.25” (6.5mm) pueden ser adecuadas dependiendo de tu técnica y del fluido que estés inyectando, pero 0.5” (12.5mm) es una apuesta segura. Más allá de 1” (25mm) es innecesariamente intimidante y doloroso sin ningún beneficio adicional.

La zona de inyección que elijas también puede jugar un papel en la longitud máxima de aguja que encontrarás aceptable. Por lo general las más cortas son preferibles, pero una longitud mayor está bien en regiones más gruesas como el muslo o los glúteos. Las inyecciones en el estómago, por otro lado, definitivamente requieren agujas muy cortas, quizá incluso más cortas que 0.5” (12.5mm), dependiendo de cómo lo estés haciendo.

\subsection{¿Importa el tamaño de la jeringa?}

\textbf{Sí, el tamaño importa.} Hay dos razones para esto. 1) Las jeringas de mayor volumen tienden a ser menos precisas, lo que lleva a una dosificación incorrecta, y 2) la física hace que las jeringas de mayor volumen sean más difíciles de inyectar. Para tener precisión en la dosificación no quieres usar una jeringa mucho más grande que el volumen que estás inyectando (es decir, para inyecciones menores a 0.1ml, usa jeringas de menos de 1ml). \textbf{Evita las jeringas de 3ml por completo si puedes.} Obviamente úsalas si es todo lo que tienes, pero realmente no están diseñadas para una tarea como esta. No me preguntes por qué los farmacéuticos parecen darlas casi exclusivamente. Una broma cruel, tal vez.

\subsection{¿Dónde compro jeringas y agujas?}

Depende de tu jurisdicción local, ya que algunas localidades prohíben la venta de agujas y jeringas a individuos como medida punitiva contra los usuarios de drogas. De no ser así, los negocios de suministros médicos y veterinarios o los minoristas autorizados de fabricantes deberían ser buenos lugares para buscar. \textbf{No se recomienda Amazon} porque no se tiene certeza sobre la calidad.

\subsection{¿Es aceptable reutilizar agujas o jeringas?}

\textbf{No. Nunca reutilices ninguna aguja o jeringa.} Ni las compartas. Para mayor claridad: Todas las jeringas, ya sea con agujas extraíbles o no, son de un solo uso. Probablemente ya lo sepas, pero solo te lo recuerdo porque ¡realmente no es bueno ni seguro hacerlo!

\subsection{¿Qué pasa si quiero hacer inyecciones pero tengo dificultades para hacerlas en mí misma?}\label{5-21}

Puede que te gusten los auto-inyectores. Como sugiere su nombre, los auto-inyectores realizan la inyección por ti. Uno como el \href{https://unionmedico.com/90-super-grip/}{\textit{UnionMedico 45/90 Super Grip}} puede tomar jeringas de 1ml lo cual puede facilitar la tarea de inyectar (pero aún así presionas el émbolo manualmente), mientras que auto-inyectores como el \href{https://www.youtube.com/watch?v=hPbhEpUN43Y}{\textit{Owen Mumford Autoject 2}} ocultan completamente la aguja de una jeringa de insulina y empujan automáticamente el émbolo hacia abajo. También hay una variedad de diseños imprimibles en 3D disponibles en línea. No he usado ninguno de estos productos y no son recomendaciones.

\subsection*{Conceptos Básicos de un Vial}
\addcontentsline{toc}{subsection}{\textemdash{} Conceptos Básicos de un Vial}

\subsection{¿Qué debería buscar al inspeccionar viales?}

Aparte de buscar signos de \textit{coring} (ver Pregunta \ref{5-23}), deberías buscar cualquier signo de decoloración, separación, contaminación, cristalización, grietas en el vidrio, fibras, pelos, etc. Un vial hecho correctamente no debería desviarse demasiado de lo habitual. \textbf{Siempre inspecciona tus viales antes de usarlos. No uses un vial que no parezca estar bien.}

\subsection{¿Qué es el \textit{"coring"}?}\label{5-23}

El \textit{coring} es cuando un pedazo de la tapa de goma se rompe y cae dentro del vial. Esto puede ocurrir con agujas de extracción demasiado gruesas, punciones repetidas en exactamente el mismo lugar, o demasiadas punciones (como en el caso de hacer inyecciones de volumen muy pequeño con un vial de volumen muy grande). \textbf{Un vial con \textit{coring} debe ser descartado inmediatamente.} También se puede usar la \href{https://www.youtube.com/watch?v=w5F0SLoMjC8}{\textit{técnica de 45-90°}} para ayudar a minimizar el \textit{coring}.

La preocupación con el \textit{coring} es que no quieres inyectarte pedazos de goma. Si hay pedazos grandes de goma, podría haber pedazos más pequeños que no puedes ver. El propósito de la tapa es proteger el contenido de los elementos, ya que un vial con un agujero en la parte superior es más propenso a la oxidación y/o al crecimiento bacteriano. \textbf{Como nota al margen: Asegúrate de quitar la tapa de metal o plástico de la parte superior de un vial nuevo.} Esto puede parecer obvio, pero algunos diseños de viales pueden ser confusos.

\subsection{¿Cuánto tiempo dura un vial?}

Un vial sellado podría durar años sin problemas si se almacena a temperaturas estables y alejado de la luz. Las preocupaciones con la edad son principalmente la oxidación del aceite portador, asumiendo que el vial fue esterilizado como debería ser. Un vial perforado que tenga un conservante en él (ver Pregunta \ref{6-17}) debería durar al menos un año o hasta que lo uses todo. La indicación de "desechar después de 28 días" en los viales de farmacia es simplemente el requisito mínimo para cuánto tiempo los fabricantes deben garantizar la esterilidad, no la vida útil máxima.

\subsection{¿Cómo debo almacenar un vial?}\label{5-25}

A temperatura ambiente estable y alejado de la luz. El calor alto y la radiación UV pueden causar degradación del aceite portador, mientras que las temperaturas bajas pueden causar cristalización. Los cristales pueden disolverse y reincorporarse, pero es una posible causa de irritación si no se disuelven completamente. Esto aplica tanto para viales sellados como para no sellados.

Adicionalmente: Siempre necesitarás reemplazar el fluido que extraes del vial con un volumen equivalente de aire para prevenir la formación a la larga de un vacío dentro del vial, lo cual puede hacer que la extracción sea significativamente más difícil. Haz esto llenando parcialmente tu jeringa con aire (hasta la misma marca que lo que planeas extraer) antes de pinchar la goma, inyectando ese aire en el vial y luego extrayendo el fluido como de costumbre. Este sencillo paso te ahorrará un dolor de cabeza en el futuro.

\subsection{¿Qué es el "espacio muerto"?}\label{5-26}

El \textit{espacio muerto} se refiere a la cantidad de fluido que se desperdicia al realizar una inyección. Este es el fluido que queda atrapado en la jeringa o en la aguja. Con una aguja/jeringa \text{luer lock} (extraíble) standard este puede ser de hasta 0.1ml, mientras que con una aguja de insulina puede ser tan bajo como 0.003ml. Reducir el espacio muerto es recomendado por razones económicas porque suma mucho estrógeno desperdiciado. \href{https://hrtcafe.net/Calc/}{Esta calculadora} puede ser útil para estimar cuánto estrógeno se desperdicia dependiendo de los suministros usados.

Una cosa a tener en cuenta: Si estás cambiando de agujas para extraer e inyectar, entonces deberías tirar del émbolo ligeramente hacia atrás antes de quitar la aguja de extracción para que el fluido dentro de la aguja de extracción no se desperdicie. Es muy menor, pero a veces puede suponer una diferencia. \textbf{Ten en cuenta también que las marcas de la jeringa miden el volumen exacto que se inyectará, así que no deberías agregar manualmente más a tu inyección para ajustar por el espacio muerto. La marca de la jeringa es precisa para la dosis.} El volumen exacto del espacio muerto depende de cada aguja en particular (es decir, una aguja de extracción y una de inyección es poco probable que tengan volúmenes de espacio muerto iguales), así que lo importante, como siempre, es la consistencia. Ver Pregunta \ref{7-7} para otra posible estrategia si te preocupa tener un espacio muerto alto.

\subsection{¿Qué hago con mis jeringas y agujas usadas?}\label{5-27}

Coloca todos los suministros de inyección usados apuntando hacia abajo en un contenedor de agujas (ya sea un contenedor dedicado de residuos biológicos o reutilizando recipientes de plástico duro como los de polvo de proteínas o detergente para la ropa). Cuando el contenedor se llene hasta tres cuartas partes, ciérralo para que no se pueda abrir accidentalmente. Etiquétalo claramente como "AGUJAS USADAS" y luego deséchalo según los requisitos de tu jurisdicción local. \textbf{Ten en cuenta que las agujas no deben ser colocadas en basura o contenedores de reciclaje.} Es probable que tu ciudad/estado/región tenga un sitio web que describa cómo y dónde desechar residuos domésticos peligrosos. Para Estados Unidos, \href{https://safeneedledisposal.org/}{puedes ir aquí.}



\section{OBTENIENDO VIALES}\label{sv}

\subsection{¿Dónde consigo viales de estrógeno para inyectar?}

En términos generales, tienes dos opciones: \textit{Proveedores farmacéuticos} y \textit{proveedores DIY} (\textit{"Do-It-Yourself"}, que es como se conoce a los proveedores \textit{no oficiales}). Los \textit{proveedores farmacéuticos} normalmente requieren una receta médica porque la TRH no está disponible sin receta (o si lo está, los viales no están incluidos) en la mayoría de los países. Los \textit{proveedores DIY} abarcan todo lo demás.

\subsection{¿Debería usar proveedores farmacéuticos o proveedores DIY?}

La elección es tuya, pero a veces no hay elección. Hay pros y contras para cada uno. Por supuesto, no hay nada que te impida obtener estrógeno de múltiples proveedores para obtener los beneficios de ambas. En muchas situaciones, puede ser recomendable.

\subsection*{Obtención de Viales de Proveedores Farmacéuticos}
\addcontentsline{toc}{subsection}{\textemdash{} Obtención de Viales de Proveedores Farmacéuticos}

\subsection{¿Cuáles son los beneficios de los proveedores farmacéuticos?}

\begin{itemize}
  \item Puedes confiar generalmente en los procesos de control de calidad y certificaciones;
  \item El seguro puede cubrirlo en parte o en su totalidad (no en todos los países);
  \item Puede ser más conveniente dependiendo de tu suerte con los médicos;
  \item El producto probablemente será consistente;
  \item \textbf{Al menos aparentar estar usando proveedores farmacéuticos puede ser requerido si estás buscando aprobación de seguro para cirugías.}
\end{itemize}

\subsection{¿Cuáles son los inconvenientes de los proveedores farmacéuticos?}

\begin{itemize}
  \item Selección reducida (o nula) de ésteres;
  \item Tiempo de espera posiblemente largo (meses o años);
  \item Puede ser necesario tener una receta (dependiendo del país);
  \item El seguro puede no cubrir los costos en parte o en su totalidad;
  \item Puede que no te lo prescriban en absoluto en tu país;
  \item Tu médico puede negarse arbitrariamente a prescribirte el medicamento;
  \item Tu médico puede negarse arbitrariamente a continuar recetándote el medicamento;
  \item Una escasez puede impedir completamente obtener el contenido de la receta;
  \item Probablemente se ciña a requisitos estrictos de WPATH o peor;
  \item Más difícil de almacenar;
  \item El acceso está sujeto a los caprichos de la situación política de tu país, lo que también significa que tu condición de persona transgénero probablemente será incluida en tu historial médico.
\end{itemize}

\subsection*{Obtención de Viales de Proveedores DIY}
\addcontentsline{toc}{subsection}{\textemdash{} Obtención de Viales de Proveedores DIY}

\subsection{¿Cuáles son los beneficios de los proveedores DIY?}

\begin{itemize}
\item Generalmente mucho más baratos en la mayoría de los lugares;
\item Disponibles en cualquier parte del mundo;
\item Obtener viales puede tomar meses o incluso años menos que las listas de espera (la única espera es el envío y la producción);
\item Fáciles de almacenar;
\item Selección completa de ésteres;
\item No hay necesidad de tratar con el sistema médico;
\item Probablemente están hechos con amor.
\end{itemize}

\subsection{¿Cuáles son los inconvenientes de los proveedores DIY?}

\begin{itemize}
  \item Probablemente no se fabrican en una sala certificadamente limpia;
  \item La calidad puede variar según el proveedor;
  \item Puede ser incómodo según el proveedor;
  \item Requiere confiar en el proveedor;
  \item Requiere encontrar un proveedor;
  \item Los proveedores DIY son más propensos a cerrar que tu farmacia local;
  \item Los tiempos de entrega de los productos pueden variar;
  \item Lo más probable es que tengas que usar criptomonedas lo cual es molesto;
  \item No puedes usar el seguro si esa era una opción para ti.
\end{itemize}
Adicionalmente, si estás buscando aprobación de un seguro para cirugías, es probable que pidan una cantidad mínima de tiempo con una receta de TRH. Esto puede ser o no una preocupación para ti.

\subsection{¿Qué tipos de estrógeno inyectable son solo DIY?}

Principalmente, el \textit{Enantato de Estradiol} y el \textit{Undecanoato de Estradiol}. Los proveedores farmacéuticos casi siempre te proporcionarán \textit{Valerato de Estradiol}, pero no siempre a una concentración de 40 mg/ml. Puede que te receten \textit{Cipionato de Estradiol} ocasionalmente, pero rara vez por encima de concentraciones de 5 mg/ml o 10 mg/ml, lo cual es molesto para dosificar. Los beneficios proporcionados por el \textit{Enantato de Estradiol} por sí solo son muy buenas razones para considerar el DIY, pero puedes obtener cualquier éster a 40 mg/ml de proveedores DIY. El \textit{Undecanoato de Estradiol} también es solo DIY para aquellas interesadas en la autoexperimentación aunque, como dije, no lo recomiendo para las inexpertas.

\subsection{¿Cuáles \textit{son} realmente los proveedores DIY?}

Los proveedores DIY pueden ser preparadores comerciales de viales, proyectos de ayuda mutua, tu amiga, ¡y tú misma si tienes un espíritu emprendedor!

\subsection{¿Dónde puedo obtener viales DIY?}

¿Qué eres? ¿Un policía? No te lo voy a decir. Ese no es el objetivo de esta guía de todas formas. Hay otros recursos que tienen esa información. Céntrate.

\subsection{¿Cómo pueden ser más baratos los proveedores DIY que los farmacéuticos?}

El costo de producir un vial es, aproximadamente, alrededor de 10 dólares, incluyendo mano de obra y costos amortizados de inversión. Esto probablemente sea una estimación alta. La mayor parte del costo de los proveedores DIY comerciales proviene de las capas de gastos operativos y de envío necesarias para mantener el anonimato. Los proveedores DIY no comerciales probablemente no tienen esos gastos fijos. Los proveedores farmacéuticos generalmente no tienen ningún incentivo para ser más baratos de lo que son.

\subsection{¿Es el DIY legal?}\label{6-11}

En la mayoría de los lugares, incluyendo América, el estrógeno no es una sustancia controlada, mientras que la testosterona puede ser o no criminalizada. Estados Unidos es una anomalía para la testosterona en este sentido, ya que otros países no criminalizan la posesión de testosterona, pero la persecución es rara de todos modos dada la amplia disponibilidad de esteroides. \textbf{Esta guía no es asesoramiento legal.}

\subsection{¿Es seguro el DIY?}

El DIY como una categoría amplia de proveedores no es ni seguro ni inseguro, pero no todas los los proveedores DIY son iguales. Cuando estamos discutiendo el tema de inyectar algo de forma segura en tu cuerpo, la verdadera pregunta es: ¿Confías en que la persona que produjo ese vial siguió adecuadamente las técnicas y procedimientos asépticos de tal manera que el vial contiene lo que quieres y nada más? Para las fuentes farmacéuticas, esa confianza es innata bajo la suposición de que existen leyes y regulaciones. Para los proveedores DIY, esa confianza debe ganarse a través de la demostración/explicación del proceso, pruebas independientes de terceros de concentración/pureza, y reputación comunitaria.

\subsection{¿Qué cosas debo buscar para saber si un proveedor DIY es confiable?}

Usa tu instinto y tu cerebro.

\begin{itemize}
  \item ¿Están abiertos a hablar contigo sobre su proceso o lo tienen listado en algún lugar? (p. ej., ¿filtran el polvo? ¡La respuesta debería ser sí!)
  \item ¿Parecen competentes con sus capacidades?
  \item ¿Han realizado tests sobre su producto?
  \item ¿Son un miembro confiable de la comunidad?
  \item ¿Han sido evaluados o respaldados por otros miembros de la comunidad en quienes confías? (es decir, inspecciones, reseñas, testimonios, etc.)
  \item Los errores ocurren pero, ¿asumen responsabilidad o intentan silenciar la negatividad?
  \item Para proveedores comerciales, ¿resuelven cualquier problema con pedidos de clientes?
  \item Para proveedores comerciales, ¿están aceptando pagos por productos que aún no se han producido sin indicar que es un pedido pendiente de producción? (¡no deberías nunca hacer un pedido de un producto sin stock!)
  \item ¿Contienen sus viales conservantes?
  \item ¿Cuánto tiempo llevan produciendo? (por buenos motivos, ¡pueden no decirlo!)
  \item ¿Cuánto producen? (por buenos motivos, ¡pueden no decirlo!)
  \item ¿Simplemente te dan \textit{mala sensación}?
\end{itemize}

Estas son solo algunas de las muchas preguntas que se pueden hacer para saber si confías en que les importa tanto como a ti la calidad de su producto.

\subsection{¿Debería aplicar standards diferentes a los distintos proveedores DIY?}

Probablemente sí. A los preparadores comerciales de viales se les debe exigir un estándar alto si les estás dando dinero a cambio de un producto ya que pueden permitirse hacerlo bien. Por otro lado, un producto de ayuda mutua que está distribuyendo viales de forma gratuita podría no ser algo sobre lo que puedas permitirte ser exigente, aunque eso no quiere decir que el producto sea probablemente mejor o peor. En cuanto a una amiga o tú misma, ¡solo tú puedes decidir eso!

\subsection*{Anatomía de un Vial}
\addcontentsline{toc}{subsection}{\textemdash{} Anatomía de un Vial}

\subsection{¿Qué debería buscar en un vial?}

Los ingredientes dentro de un vial pueden categorizarse como \textit{"activos"} y \textit{"excipientes"}. El \textit{activo} es el éster de estrógeno en nuestro caso, y los \textit{excipientes} son todo lo demás. Generalmente hay tres o cuatro ingredientes: 1) El éster de estrógeno; 2) el aceite portador; 3) el conservante; y, opcionalmente, 4) cosolvente(s). Ya hemos cubierto los ésteres de estrógeno en la Sección \ref{td} "TIPOS Y DOSIS". Los viales farmacéuticos casi siempre tienen los cuatro ingredientes.

\subsection{¿Qué aceite de transporte debería buscar en un vial?}\label{6-16}

Esta es una cuestión de preferencia, tolerancia personal y posiblemente alergias. \textbf{La variable principal relevante para ti es la viscosidad porque eso afecta a la comodidad y conveniencia de la inyección.} Como se discutió anteriormente (ver Pregunta \ref{5-16}), los aceites menos densos pueden ser usados con agujas de mayor calibre de manera más cómoda sin dificultad al extraer e inyectar. \textbf{Los aceites portadores más comúnmente usados para TRH son el aceite de ricino (\textit{"castor oil"} en inglés)y el aceite MCT.} El aceite de ricino es el más grueso de los aceites comúnmente usados, pero también tiende a resultar en menor cantidad de irritación, por lo que los viales farmacéuticos típicamente lo usan. El aceite MCT es el aceite más ligero comúnmente usado, pero algunas personas lo encuentran más irritante que otros aceites y solo está disponible en DIY. El aceite de semilla de algodón y el aceite de pepita de uva se usan ocasionalmente, pero generalmente no por los fabricantes de HRT. Otros aceites como el de girasol o sésamo o cualquier otro ocasionalmente encuentran uso pero generalmente no se recomiendan. Dependiendo de tus circunstancias, esta pregunta podría no importarte, podrías no tener una opción, o podría ser un requisito estricto.

\subsection{¿Qué conservantes debería buscar en un vial?}\label{6-17}

El conservante más comúnmente usado en los viales inyectables es el \textit{alcohol bencílico} (BA, del inglés \textit{benzyl alcohol}) en baja concentración. Esto es obligatorio y no está sujeto a debate. \textbf{Nunca deberías usar un vial sin conservante.} Para personas que sean (raramente) alérgicas, el \textit{chlorobutanol} es un conservante alternativo comúnmente usado, pero casi nunca por proveedores DIY, así que requeriría buscar fórmulas farmacéuticas específicas.

\subsection{¿Qué cosolventes debería buscar en un vial?}

El cosolvente principalmente usado es el \textit{benzoato de bencilo} (BB, del inglés \textit{benzyl benzoate}) que reduce la viscosidad de la solución resultante. Esto es técnicamente opcional, pero generalmente se recomienda para consistencia del lote y en muchos casos es necesario dependiendo del aceite portador (sobre todo en aceite de ricino) y la concentración deseada. Algunas personas lo encuentran irritante, otras no.


\section{RESOLUCIÓN DE PROBLEMAS}

\subsection*{Incertidumbres con la Dosis}
\addcontentsline{toc}{subsection}{\textemdash{} Incertidumbres con la Dosis}

\subsection{Mis niveles no son los que esperaba. ¿Por qué no?}

Hay una serie de posibilidades. Recuerda primero que las estimaciones del modelo no pueden tener en cuenta toda la gran cantidad de factores que pueden causar alguna desviación. Recuerda también que se necesitan múltiples inyecciones hasta que alcances la estabilidad, así que si acabas de cambiar tu dosis, esa puede ser la razón. Verifica por cuatruplicado con un amigo que estás inyectando lo que crees que estás inyectando. Eso es problema más común de lo que podrías pensar, pero con los proveedores DIY también es posible que la concentración sea menor de lo anunciada debido a la falta de experiencia o a equipos menos precisos. En ese caso, al inyectar es posible que simplemente necesites inyectar un poco más cuando uses ese vial. \textbf{Pero recuerda, lo más importante es cómo te sientes, no tus niveles.} Ten en cuenta que incluso las farmacias de compuestos profesionales pueden producir viales defectuosos no detectados por el control de calidad, ¡por más raro que eso sea, afortunadamente!

\subsection{¿Puedo comparar los niveles de diferentes análisis si no me los hice en mi nivel mínimo?}

\textbf{No.} No con precision, al menos. Esto es parte de por qué siempre deberías hacer los análisis en tu \textit{nadir} (punto más bajo, \textit{trough} en inglés). Es decir, horas antes de tu momento normal para tu próxima inyección; eso es lo que quieres. Eliminar tantas variables como sea posible hace que los datos sean mucho más útiles para ti. Si no hay nada más que saques de esta guía, por favor, solo haz tus análisis en tu nadir/nivel mínimo.

\subsection{Me siento realmente mal en mis días de nadir. ¿Qué debería hacer?}\label{7-3}

En la mayoría de los casos, o bien la dosis es demasiado baja o la frecuencia es demasiado baja. Esto es de vital importancia para el \textit{Valerato de Estradiol} y el \textit{Cipionato de Estradiol}. Ajusta tu dosis dentro del rango listado o ajusta la frecuencia. Encuentra lo que funciona para ti. También es posible en particular con el \textit{Valerato de Estradiol} que tu dosis pueda ser demasiado *alta* en lugar de demasiado baja, ya que la alta variabilidad de niveles a lo largo de tu ciclo puede ser la causa de esta sensación de bajón. En resumen: Cambia a \textit{Enantato de Estradiol} si puedes.

\subsection*{Dificultades con las Inyecciones}
\addcontentsline{toc}{subsection}{\textemdash{} Dificultades con las Inyecciones}

\subsection{La inyección es más difícil de hacer cuando está fría. ¿Qué debería hacer?}

Calienta el vial antes de extraer del mismo, luego calienta la jeringa antes de inyectar. Frotar el tubo de la jeringa entre las manos debería ser más que suficiente para calentar el líquido. Hacer que esto se convierta en un hábito todo el tiempo debería mejorar tu consistencia en las inyecciones.

\subsection{La inyección duele más cuando está fría. ¿Qué debería hacer?}

Calienta tu sitio de inyección antes de inyectar. Relajar los músculos con un masaje o una ducha caliente (específicamente, aumentar la temperatura con el agua apuntada hacia tu sitio de inyección antes de que salgas) antes de inyectar puede ayudar.

\subsection{Sangré después de mi inyección. ¿Me moriré?}

No. Esto significa que lo más probable es que simplemente hayas tocado una vena o un capilar, lo cual puede suceder a veces. Es posible que te salga un pequeño hematoma o que sientas un poco más de dolor. Ponerte una tirita mona hará que sane más rápido.

\subsection{Había algo de aire en mi jeringa. ¿Me moriré?}\label{7-7}

No. Aunque obviamente no quieres inyectar solo aire y puede afectar a la dosis si hay demasiado en la jeringa, una pequeña cantidad de aire por debajo de 0.1ml casi con certeza no va a causarte problemas. De hecho, podría ser recomendado en algunos casos. Por ejemplo, la \textit{técnica air lock} (bloqueo de aire, una técnica estándar para inyectar fluidos que son irritantes o pueden manchar; no es un conocimiento crucial para HRT) generalmente implica inyectar 0.1-0.3ml de aire, así que no tienes nada de qué preocuparte. No estás haciendo inyecciones intravenosas.

\subsection{Algo del fluido se salió. ¿Se desperdició mi inyección? ¿Me moriré?}

No. Las fugas pueden ocurrir por toda una variedad de razones y rara vez pueden ser suficientes como para suponer diferencia, así que no necesitas hacer otra inyección. Para el futuro, asegúrate de dejar la aguja dentro durante 5-10 segundos antes de retraerla y luego aplicar presión después. Podrías considerar usar la técnica de bloqueo de aire mencionada anteriormente o \href{https://www.nurse.com/nursing-resources/definitions/what-is-z-track-method/}{el método de rastreo en Z} si estás particularmente preocupada por las fugas.

\subsection{A veces me siento muy dolorida después de una inyección. ¿Me moriré?}

No. Asumiendo que por lo demás has seguido todas las sugerencias de esta guía, a veces, por una razón u otra, el depósito de líquido termina situándose en un lugar incómodo. Más suerte la próxima vez. \textbf{¡Asegúrate de alternar los lugares de inyección!} No quieres que se acumule tejido cicatricial a largo plazo, y si una zona ya está dolorida, no quieres hacerla más dolorida.

\subsection{Estoy experimentando mucha picazón e irritación después de inyectarme. ¿Me moriré?}

Probablemente no. Hay una serie de causas posibles. La infección es la causa más preocupante, pero es poco probable en la mayoría de los casos. \textbf{Ve inmediatamente al médico si estás experimentando fiebre, dolor severo, dolores musculares, pus, rayas rojizas u otros signos de infección.} En la mayoría de los casos, sin embargo, la irritación como picazón, enrojecimiento, hinchazón ligera, calor, etc., es el resultado del uso de un vial cuyo estrógeno y aceite se han separado (cristalizado o salido de la solución). Mira abajo. Es posible que estés teniendo una reacción al aceite portador, pero si repentinamente estás experimentando problemas después de algunas inyecciones sin ningún problema previo, es muy probable que el contenido del vial esté cristalizado.

\subsection{Mi vial tiene cristales dentro. ¿Aún puedo usarlo?}

Probablemente significa que tu vial está demasiado frío. Caliéntalo y agítalo suavemente para reincorporarlo. Si los cristales no desaparecen, entonces es posible que el contenido del vial se haya separado completamente. Con mucho más calor y agitación, los cristales podrían reincorporarse, pero es más simple y seguro reemplazar el vial si puedes.



\section{PROGESTERONE}

\subsection{Do I want to take progesterone?}

\textbf{Probably.} This is a controversial question for some reason. Detractors (namely, doctors) will argue that there’s no studies to show that it plays a role in feminization therefore it should not be taken. Aside from transfeminine subjects being woefully understudied, heuristically speaking, progesterone is a key female sex hormone that plays an important role in the brain and in many functions throughout the body. Regardless of physical feminization, it is an important hormone for good health that should not be lightly overlooked.

\subsection{What is the difference between “progesterone” vs “progestin” / ”progestogen”?}

The class of hormones, both natural and synthetic, that activate the progesterone receptor are “proges\textbf{togens}”. The natural, bioidentical, and most important progestogen is “proges\textbf{terone}”. Synthetic progestogens are “proges\textbf{tins}”. These three terms are mistakenly used interchangeably in scientific literature and in clinical settings, likely causing much of the broader confusion regarding the role of progesterone in HRT, despite the fact that they are \textbf{not} equivalent.

\subsection{Do I want progesterone or a progestin?}

Progesterone. You want bioidentical progesterone.

\subsection{What’s wrong with progestins?}

Progestins, most typically \textit{medroxyprogesterone}, \textit{medroxyprogesterone acetate}, or \textit{levonorgestrel}, are generally associated with the negative side effects and long term risks (breast cancer, blood clots, depression, etc) that are falsely attributed to progesterone. They are not bioidentical which means they do not behave the same as progesterone and thus cannot be directly compared.

\subsection{What does progesterone do for feminization?}

Progesterone is believed to play a role in breast development and libido in particular, but as mentioned it’s a key hormone aside from its outward appearance effects. It does also have some antigonadotropic (i.e., it contributes to testosterone suppression) properties which can be sometimes relevant.

\subsection{Does it matter when I start progesterone?}

It is unknown. There is some belief that starting too early may harm breast development long term, but this is purely theoretical and contrary anecdotal evidence makes the answer unclear. The conservative estimate is waiting roughly a year into HRT (until Tanner Stage 3 or 4) in the possible chance that it does matter.

\subsection{How is progesterone normally taken?}

Aside from topical applications, the main form is via a pill. It is prescribed as an oral pill but is most effective when taken as a suppository. Topical sprays and creams can also work very well.

\subsection{Are you serious that progesterone should be taken as a suppository?}

Progesterone metabolizes entirely differently when taken orally vs rectally due to passing through the liver when taken orally. Oral progesterone primarily converts to \textit{allopregnanolone} which can cause heavy drowsiness, whereas rectal progesterone primarily converts to progesterone itself which is what we want (although some still converts). Some people take additional oral progesterone as a sleep aid, but please note that too much \textit{allopregnanolone} can sometimes lead to negative mental health side effects.

\subsection{How do I take progesterone as a suppository?}

Just a bit of water on the pill should work, then dry off and wash your hands. Obviously, don’t go to the bathroom for the next hour or so, so doing it before bed is best. If you are having issues with it not dissolving then you can try piercing the capsule but usually should be no issue. Be aware that if you use large homebrew suppositories made using coconut oil, the large volume of coconut oil will not want to stay in you.

\subsection{How much progesterone should I take?}

For pills, Standard dosage is 100-200mg daily at night. It is a rather arbitrary dosage; 200mg is the max that most doctors will prescribe. Some people take more than 200mg on occasion, but be aware that spiking your levels may lead to an unpleasant crash. See question below.

For topical applications, nobody can tell you with certainty due to the high variability of the delivery medium, nor is there any clear guidance on desired levels, or even frequency (likely daily), as progesterone is simply understudied. Because of this, I would advise titrating your dosage so that you understand how progesterone affects you.

\subsection{Is there any benefit to “cycling” progesterone?}\label{8-11}

No. Some people do this to mimic a cis woman’s menstrual cycle, but there is no reason to believe there is any benefit to this and it may cause negative PMS symptoms. The only exception is if you have good reason to suspect that you have an intersex condition involving a uterus that you are managing. I discourage it otherwise. See Question \ref{11-10}.

\subsection{How long should I take progesterone for?}

For as long as you plan to take estrogen and for as long as you want to. So, probably forever.

Sometimes people (or doctors) arbitrarily say to only take progesterone for X years. There is zero theoretical or empirical reason to suggest that this is sound advice. It's about as coherent as if someone (or a doctor) asked how long a trans person planned to take HRT for\textemdash{}oh wait never mind they do ask that.

\subsection{Can progesterone convert into \textit{dihydrotestosterone} (DHT)?}

No. Well, strictly speaking yes, but also no. It is largely a myth, although \href{https://whsah.co/posts/rethinking-progesterone-and-androgens/}{as outlined in detail by alix [sic] in this article}, for cases of people with \textit{nonclassical congenital adrenal hyperplasia} (ncCAH) progesterone can cause some negative side effects of increased androgenic activity. In those cases, discontinuing progesterone is recommended along with seeking out a formal diagnosis/treatment for potential adrenal disorders.

\subsection{Is there any benefit to topical progesterone applications in addition to pills?}

Maybe. It's a possible alternative to pills, especially in the case of someone with a peanut allergy since the most common pill manufacturer uses peanut oil, but again dosage is unclear. Some people find more progesterone fun, if nothing else. Be safe and have fun.

For clarity: Apply creams to your inner thigh region (elsewhere if directed), or optionally on scrotal skin (it's thin and highly vascular) in the case of sprays. And no, applying progesterone to your breasts directly is unlikely to make them grow bigger or faster compared to otherwise.

\subsection{Can I snort progesterone powder?}

Please don’t. It’s hell on your sinuses. It isn’t hard to make your own topical progesterone spray and there are guides out there. Do that instead. It’s significantly more effective, consistent, and safer.

\subsection{Where can I get progesterone?}

Progesterone tends to be more expensive through DIY sources due to the higher mass of hormones required, so ideally get it through pharmaceutical sources covered by insurance. There is also the option of grey market foreign pharmacies, which are simply pharmacies in another country, although these often require some hurdles to purchase from. Topical progesterone creams are available OTC in some locations, although it is not always the most economical depending on the concentration.

\subsection{I would like to read more about progesterone in an HRT context. What resources should I read?}\label{8-17}

Originally I linked a document here but I opted to remove it due to a number of faults that can be misleading. The problem with progesterone is that literally nobody agrees about a single aspect of it. I don't know a single source or study that people agree is good. Hell, people don't even agree if the word starts with the letter “P”. \textbf{The crucial thing to know is that progesterone is not strictly required for full feminization or good breast development, but assuming that it's not contraindicated for you, it's probably worth taking.}

It should be noted that for the entire category of progestogens there are countless myths and falsehoods invented whole cloth by both proponents and detractors alike which does not make discerning truth from the already-sloppy scholarship any easier. Fantastical claims of magical benefits and fearmongering of alleged risks based on nothing are both equally unhelpful, although the latter is worse in my opinion when comes from a medical authority, whether neglectful or malicious.

\subsection{Does progesterone interact with any other drugs related to HRT?}\label{8-18}

If you are taking 5$\alpha$-Reductase Inhibitors like \textit{finasteride} and \textit{dutasteride} (see Section \ref{AA} “ANTIANDROGENS”, or keep reading), these can affect how progesterone naturally breaks down into \textit{allopregnanolone} which can cause adverse mood effects in some people, irrespective of how you are taking progesterone. It is not fully clear how much the administration route for the 5$\alpha$-Reductase Inhibitors (i.e., topical vs oral) makes a difference, but lower systemic absorption via topical application may mitigate these side effects. It is recommended to not take either of those if you are someone affected by this interaction, but it is not in all cases anyway. Note that these depressive effects may be felt for up to a month after stopping.


\section{TESTOSTERONE}\label{T}

\subsection{Why don’t we want zero testosterone?}

Testosterone is an essential sex hormone which plays a key role in your health and well-being. We want to suppress it for feminization, but near-zero testosterone (less than 10 ng/dl, or 0.35 nmol/L) can cause issues such as poor libido, low energy, low strength (fatigue beyond just the strength loss of HRT), poor concentration, trouble sleeping, etc. Notably, issues very similar to having too little estrogen. Cis women also have more than zero testosterone, so that need not be the fear. \textbf{Adequate hormone levels are important!}

\subsection{Are there ever cases where I would want to supplement testosterone?}\label{9-2}

Yes. If you are experiencing the issues of the above and your estrogen levels are otherwise good, it’s possible that you might want to supplement with a microdose of testosterone. If you wanted to improve your erectile function, minimize any atrophy before bottom surgery, or otherwise wanted to experiment with your hormones to see what feels best for you, then that might be a reason to explore testosterone in a different context that you can hopefully appreciate more compared to pre-HRT.

\subsection{If I wanted to supplement testosterone, how would I do it?}

There’s a few possibilities. Testosterone comes in either injections or topical gels/creams, similar to estrogen as already discussed. Topical is more likely what you are going to be prescribed. Topical applications have the downsides that we have discussed for estrogen, but those are less of a concern here when precise levels are less important.

\subsection{What are the topical forms of testosterone?}

There is gel and cream. Gel is typically what will be prescribed, but some compounding pharmacies are able to make low-penetrating cream if someone wanted just topical application on the genitals. The latter is harder to get and generally more expensive, however.

\subsection{Does it matter where I apply the testosterone?}

It depends on if you have gel or cream. If you have the kind of localized cream as mentioned above, you would apply it as directly as mentioned. Otherwise, shoulders and upper arms are where gel should go. Make sure not to touch things until long after it dries!

\subsection{How much and how often should I apply testosterone?}

Season to taste. This largely depends on how you are feeling. If you have too much, you might start to experience side effects of testosterone (e.g., oily skin and body hair), but only you can say what is preferred for you. A weekly injection of 5-10mg of \textit{testosterone cypionate} might work for you, but in the case of 1\% topical gels which are often disbursed in 25/50mg packets, there is more variability. You almost never want even half a packet, and definitely not daily. I would suggest starting with much less than you think to see how you feel.

\subsection{Where would I get testosterone?}

If you are an American, you would have to get a prescription or ask any juicer at your closest Planet Fitness. Elsewhere, it depends on what gym chain is closest to you. Disclaimer: This is a joke. See Question \ref{6-11}.

\subsection{Are other steroids equivalent to testosterone in an HRT context?}

Anabolic-androgenic steroids, i.e., drugs that are structurally similar to testosterone, are not all equivalent. Commonly used black market steroids like \textit{trenbolone acetate} have a laundry list of undesirable side effects, but steroids like \textit{nandrolone decanoate} are occasionally used for postmenopausal cis women due to their relatively low androgenic properties which make them very favorable for transfeminine individuals. Regardless, in America it is unlikely you will be prescribed anything other than testosterone itself, if you are able to get a prescription at all.

\subsection{What is the relationship between testosterone and \textit{dihydrotestosterone} (DHT)?}

\textit{Dihydrotestosterone} is primarily synthesized from testosterone via the 5$\alpha$-Reductase enzyme with around 5\% of testosterone in your body being converted. Generally speaking, if testosterone levels are suppressed (or if you have had bottom surgery) then there should not be much left to convert, but systemic levels won’t be zero because it is still locally produced. Depending on your body, this would be the main reason that you might want to consider supplementing with a 5$\alpha$-Reductase Inhibitor antiandrogen as discussed in the following section. As a reminder, \textit{dihydrotestosterone} is the hormone that is responsible for body hair and hair loss.

\textbf{For any trans mascs reading this,} I will make a brief detour to note that at time of writing it is not clear what role the hormone plays with bottom growth regarding speed or total size as it relates to 5$\alpha$-Reductase inhibition. That is to say: it is known that \textit{dihydrotestosterone} plays a primary role in penile development, but it’s not clear how directly the lack thereof affects a trans masc person. Applying knowledge of micropenis treatment, we know that a topical cream is more effective than exogenous injections particularly with how \textit{dihydrotestosterone} cream is useful when a patient doesn’t respond to testosterone (particularly in the case of 5$\alpha$-Reductase deficiencies). So, food for thought. Someone get Oliver Longdick to handle the rest of this.



\section{ANTIANDROGENS}\label{AA}

\subsection{What are “antiandrogens”?}

\textit{Antiandrogens}, commonly also referred to as “T blockers” or just “blockers”, as the name(s) may suggest prevent androgens (that’s what testosterone is) from acting on your body. There are many types of antiandrogens and they are commonly prescribed as part of an HRT regimen. They are needed if someone still produces testosterone and is not doing a form of HRT conducive to monotherapy, such as injections, but they are usually not desirable. It also should be noted that (most) antiandrogens do not reduce testosterone levels in any way that matters but instead simply reduce/negate effects on the body. This is relevant when interpreting lab results and such.

\subsection{Why wouldn’t I want antiandrogens?}

The main issue with most antiandrogens is that they generally have very undesirable side effects that are superfluous if testosterone is suppressed in the first place by having enough estrogen, so those side effects are being experienced despite\textemdash{}in most cases, at least\textemdash{}being rendered unnecessary by a reasonably-dosed monotherapy regimen. Bottom surgery of any kind also makes antiandrogens unnecessary in most cases.

\subsection{When might I want antiandrogens?}

If you are not most cases, if you desire peace of mind, or if your insurance requires a prescription on file before they will cover a procedure, then you may want antiandrogens. The medications used as antiandrogens might have other effects that may be desirable outside of their antiandrogen properties depending on your health situation. Additionally, if you are supplementing androgens, you may want a \textit{dihydrotestosterone} (DHT) blocker to minimize side effects related to body hair and hair loss, but be aware that this may not be the case if you are not using bioidentical testosterone (e.g. \textit{nandrolone decanoate}) because not all androgens behave the same.

\textbf{It should be noted that temporarily using antiandrogens at the start of HRT when planning to perform monotherapy is unlikely to be necesssary nor is it recommended.} There is an adjustment period that you will experience regardless while your body adapts to the change in your primary hormones, so there is no need to overcomplicate what you are doing. Don't worry about it.

\subsection{What kinds of antiandrogens are there?}

The main medications taken as general testosterone blockers in an HRT context are \textit{spironolactone}, \textit{bicalutamide}, and \textit{cyproterone acetate}. The main medications taken to block the conversion of testosterone into \textit{dihydrotestosterone} (DHT) called “5$\alpha$-Reductase Inhibitors” (5-ARI) are \textit{finasteride} and \textit{dutasteride}. There are also GnRH agonists like \textit{leuprolide} and \textit{triptorelin}, but both of those are more often used as puberty blockers in minors, although in parts of Europe they are used for adults as well.

\subsection{When might I want to take \textit{spironolactone}?}

Due to the heroic dosages and significant negative side effects required for it to function as an antiandrogen in most cases, the only time I would ever recommend taking \textit{spironolactone} would be if you would benefit from its other effects such as its antimineralocorticoid (i.e., blocking \textit{aldosterone}) properties as it relates to blood pressure management or edema. \textbf{If you insist on taking \textit{spironolactone}, please do not take more than 100mg daily.} It has a bad reputation for a reason. “The Devil”, as it were.

In case you are unfamiliar, some of the many side effects include: brain fog, lethargy, poor memory, increased urination frequency, low blood pressure, low sodium / electrolyte imbalance, etc. In other words, \textit{spironolactone} is a blood pressure lowering diuretic that is a mediocre antiandrogen which is typically prescribed at high dosages in an otherwise-healthy population for questionably-effective off-label use. In any other healthcare context this would (or SHOULD!) be highly unadvisable given the undesirable side effect profile and the widely-available preferable alternatives that already exist, but that's the state of trans healthcare for you.

\subsection{When might I want to take \textit{bicalutamide}?}

If you are going to take an antiandrogen, \textit{bicalutamide} is likely the one to take. It is generally well tolerated, barring 1\% cases of abnormal liver function test results and symptoms of liver dysfunction, but otherwise performs the job with relatively minimal side effects. \textbf{If you take \textit{bicalutamide}, ensure regular liver function tests to make sure that your results are in range.} The liver risks are dependent on your body rather than cumulative so any problem would likely present itself within the first year. Otherwise, there should be no issues.

\subsection{When might I want to take \textit{cyproterone acetate}?}

Likely never. Take \textit{bicalutamide} instead.

The long term risk profile is poor and there is no situation that I can think of in which I would recommend this over an alternative solution. You can do everything \textit{cyproterone acetate} can by just taking more estrogen and adding progesterone to your regimen.

\subsection{When might I want to take \textit{dutasteride}?}

If you are extremely concerned about possible hair loss and/or want to maximize your chances for hair regrowth, you may want to take \textit{dutasteride}. If your testosterone is otherwise suppressed then it theoretically shouldn’t have much benefit as your \textit{dihydrotestosterone} levels should be relatively low, but bodies can be complicated, so it may be something of interest to you. Also, see Question \ref{11-14}.

It should be noted that \textit{dutasteride} can cause adverse mood effects in some people, in which case stopping is strongly recommended. Note as well that these depressive effects may be felt for up to a month after stopping.

\subsection{When might I want to take \textit{finasteride}?}

If \textit{dutasteride} is not something prescribed to you or if your insurance mandates \textit{finasteride} specifically to cover a hair treatment. Otherwise, \textit{dutasteride} is preferred as it is more effective and better tolerated.

It should be noted that \textit{finasteride} can cause adverse mood effects in some people, in which case stopping is strongly recommended. Note as well that these depressive effects may be felt for up to a month after stopping.

\subsection{Where can I get antiandrogens?}

Aside from being prescribed them by your doctor or perhaps available over-the-counter, there is also the option of grey market foreign pharmacies. These are simply pharmacies in another country, although these often take some hurdles to purchase from. \textit{Dutasteride} and \textit{finasteride} are generally the easiest to get over-the-counter because of their commonality as hair loss medication.



\section{MYTHS AND MISCS}\label{MM}

\subsection*{Common Questions}
\addcontentsline{toc}{subsection}{\textemdash{} Common Questions}

\subsection{Should I be worried about blood clots?}\label{11-1}

Yes and no. It is true that there is a correlation between estrogen dosages/levels and blood clot risk, but this is primarily related to the route of administration and the type of estrogen. Synthetic estrogens are the rightful cause of scorn and do lead to significantly increased blood clot risk, but bioidentical estrogens are not as concerning. In particular, the route of administration makes a major difference. Oral bioidentical estrogen passes through the liver which is what causes the increased blood clot risk. Injections bypass the liver, and there's no evidence to suggest nor reason to believe that injections of bioidentical estrogen provide any significant risk increase beyond the innate differences between testosterone and estrogen. The pervasive fearmongering towards all estrogen has persisted for decades despite these differences.

\textbf{If you are undergoing surgery, please know that pausing hormones out of concern for blood clots is no longer recommended by WPATH.} Many surgeons still include it in their pre-surgery guidelines out of concern for blood clots, but this is torture that has been disproven and even WPATH doesn't recommend it anymore. Remarkable, I know. Per \href{https://www.tandfonline.com/doi/pdf/10.1080/26895269.2022.2100644}{WPATH SOC 8 Statement 12.20}: \blockquote{After careful examination, investigators have found no perioperative increase in the rate of VTE [KT: \textit{venous thromboembolism}, i.e. a blood clot] among transgender individuals undergoing surgery, while being maintained on sex steroid treatment throughout when compared with that among patients whose sex steroid treatment was discontinued preoperatively (Gaither et al., 2018; Hembree et al., 2009; Kozato et al., 2021; Prince \& Safer, 2020).} I should put this in another question entirely, but to not break links, it would have to be at the bottom of a section and I think this is too important for that, so I note it here. A very important clarification that I should have had sooner.

\subsection{Is it okay to use nicotine while on HRT?}\label{11-2}

This is related to the above question. \textbf{Nicotine usage on HRT, especially if you’re on pills, compounds your risk of a blood clot on top of all the other reasons that nicotine is not good.} This extends to all forms of nicotine usage, but obviously smoking is by far the worst. You really do not want a blood clot. Even if you are not on pills, nicotine disrupts the way estrogen is metabolized and can lead to significantly reduced feminization effects. This aspect is understudied but community anecdotal reports are common. It’s not easy to quit, but I believe in you. There are good resources out there and strategies like tapering down by using patches really do work. You got this.

However, to be abundantly clear, \textbf{this does not mean that you cannot or should not take estrogen. The downsides of not taking estrogen at all far exceed the downsides of using nicotine.} This section is simply seeking to make you aware of any increased risks and potentially slower transition as a very strong recommendation and encouragement to quit. One step at a time.

\subsection{Is there benefit to starting at a low dosage vs a high dosage?}

To the best of knowledge, no. Sex hormones are not like other drugs that need to be titrated to manage side effects as we know the dosages that work for the majority of people, so personally I view “starter dosages” and “antiandrogen first” regimens as medical abuse. Some people believe that mimicking the slow timeline of puberty might be best (even though there are far more things happening than just estrogen levels), but there’s no evidence to support this. An orchiectomy day one might be best for all we know, but who is going to do that the moment they decide they are trans and/or want to start HRT?

Reframing this in another way: \textbf{there is no reason to believe that “starting slowly” on a dosage below the typical range is advantageous or preferable for feminization outcomes.} There isn't a concern of going “too fast” or anything like that. Both doctors and other trans women seemingly invent new myths by the day.

\subsection{Does body weight affect dosage?}

No. Because there is no “optimal” blood level for estrogen and because the therapeutic range of acceptable levels is so wide, \textbf{body weight does not meaningfully affect dosage for HRT.} It is for the same reason that slight deviations in dosage are unlikely to affect how you feel. There is no such thing as being “too light” or “too heavy” for HRT in any capacity.

\textbf{For the same reason, adjusting your dosage in increments of 0.1mg (note that says milligrams, not milliliters) is a difference that should not be expected to be perceived simply because our bodies are not sensitive enough to such exact measurements,} let alone the high possibility of imprecision when performing an injection that makes that certainty of this measurement unlikely. In other words, the accuracy of your dosage is more important than the precision.

\subsection{Is there such a thing as starting estrogen too late?}\label{11-5}

\textbf{No.} No matter when you start, estrogen is able to do a LOT and a proper regimen will be able to have powerful results. Sex hormones are some of the strongest hormones in our body in terms of our appearance. Everybody always wishes that they could’ve started sooner, but that’s no reason not to start now. Even if you’ve been on estrogen for years, there is still a benefit to be had in improving the quality of your regimen.

\subsection{Does feminization / breast development stop after X years?}\label{11-6}

\textbf{No.} There is not an arbitrary time where estrogen suddenly stops working. Various numbers are given and usually it’s either 1) entirely made up or 2) pointing to a study that only went for X years. Doctors in particular love to tell trans women not to expect more than B cup breasts (which isn’t even how breast sizing \textit{works}, but I digress) or for any growth after 2 years, but this is simply not true. There are cases of people who restarted estrogen after stopping for many years and still experienced new growth.

\subsection{I haven’t seen any changes in years on injections. Would swapping back to pills make a difference?}

Maybe, but maybe not. There are some anecdotes of people swapping back from injections to pills (or adding pills on top of injections) and experiencing more breast growth after “stalling out”, but the mechanism is not clear. There is speculation that the E1:E2 ratio (\textit{estrone} : \textit{estradiol}) heavily weighted towards E1 with oral pills compared to E2 for injections might make a difference for some people, although \textit{estrone} is not typically associated with feminization. There likely are other factors at play, but you are free to experiment if you wish. Data is limited.

\subsection{Is low energy and low libido normal on HRT?}

Generally, no. How libido is expressed changes in the beginning, but the vast majority of the time that someone experiences abnormally low libido it’s because they haven’t gotten their hormones sorted. The same goes for low energy. Get your hormones squared away, and barring that, check your diet/vitamins next. Make sure you don’t randomly have critically low vitamin D levels or something like that. It happens more often than you think.

\subsection{I hear about [random drug / strategy] that my friend said helps feminization. Does it actually?}

Maybe, but probably not. There is a lot of wild speculation about ways to achieve feminization goals, but many of them are akin to snake oil or have potentially serious risks far beyond HRT itself. You have the right to bodily autonomy and I cannot stop you, but I can encourage you to be smart about what you are doing. The more you get into the weeds of biology as it relates to transition, the shakier the ground becomes as quality data becomes less and less available. Desperation can lead to a lot of unwise and dangerous decisions. So be smart, and be safe.

\subsection{Do we want to mimic the estrogen cycle of cis women?}\label{11-10}

Arguably no. This is controversial, but I am of the belief that because we (well, most of us) do not have a uterus and corresponding menstrual cycle synced to our hormone levels, then there is no reason we should strive to copy that behavior. This is an \textit{is-ought} problem, in my view. The primary hormonal concern for most trans women is testosterone suppression which necessitates consistently high enough levels (barring post bottom surgery, where there is no testosterone to suppress), so high fluctuation and/or relatively low levels are likely to cause undue distress. You’re welcome to experiment, of course. Especially if testosterone suppression is no longer a concern for you. See Question \ref{8-11}, and see below.

\subsection{Do trans women experience periods?}

Similar to the last question, it’s important to understand what is happening. The unique hormone curve produced by your particular ester, your dosage, and your frequency can cause changes in your mood as your estrogen levels oscillate between injections. Some trans women liken this phenomenon to a period, but the underlying cause for these physiological changes is different and is usually a sign that your regimen needs tweaking so that you feel the best that you can as suffering is not virtuously feminine. Pain and discomfort are not requirements for womanhood nor should we assert ourselves based on bioessentialist arguments. The exception here are the intersex trans women who have a uterus and literally are having a period, in which case: yeah duh. See Question \ref{11-35}.

\subsection{Can too much estrogen convert to testosterone?}

\textbf{No.} Aromatase is the enzyme responsible for converting testosterone into estrogen, but there is no mechanism to convert estrogen into testosterone. This cannot happen. This is a completely false myth and you should be immediately wary of the knowledge level of anyone who says it to you. Unfortunately, it is doctors who repeat this myth the most.

\subsection{Does bottom surgery cause an increase in testosterone?}

No. This is not a thing. There is not a magic mechanism that suddenly causes testosterone to increase the moment that testicles are removed. Even if magic was stored in the balls, this simply isn’t how hormone production works. “Well, your adrenals…” They don’t work like that either. The only possible rare exception would be undiagnosed adrenal hyperandrogenism conditions that were suppressed by an antiandrogen like \textit{spironolactone} prior to surgery which might show itself after antiandrogens are ceased. Please stop repeating this myth.

\subsection{How do I prevent/revert hair loss?}\label{11-14}

Mechanically, it is pretty simple. A standard HRT regimen alone is borderline magic (don’t ask where the magic is stored) in this regard already, but the inclusion of 5$\alpha$-Reductase Inhibitors (5-ARI) as discussed in Section \ref{AA} “ANTIANDROGENS” is recommended in more extreme cases to completely halt any loss. Topical minoxodil 5\% is the only thing that works to firm up your hairline beyond hormones alone, but keep in mind that aside from miracle cases, you’re only saving dying/dormant follicles. Dead follicles don’t come back.

If this alone is insufficient for you, hair transplant technology has improved significantly. The Follicular Unit Extraction (FUE) procedure is what you want to look into. Here is where in the future I will link a guide written by an expert on getting insurance to cover that, once she writes it. This is peer pressure. Watch this space.

\subsection{Does exercise affect feminization?}

Probably. HRT causes gradual body recomposition, so you can help encourage your body to shift through exercise. Keep in mind that this process is VERY SLOW, so it is crucial that you eat enough to fuel how patient you have to be. The growth hormones from muscle stimulation via strength training also play a role in breast development, so it’s probably a good thing even aside from the rest of the obvious health benefits of exercise.

This is NOT just the writer’s barely-disguised fetish; strength training is important for your health! I mention this because a lot of trans women believe that touching a dumbbell will make them look like the hulk. I get it, but if you have no testosterone in you and you aren’t on steroids, then you aren’t going to look like that. Let alone the constant time, effort, and diligence required to even get close.

\subsection{What should I exercise then?}\label{11-16}

Cardio is useful for living which is important. Lower body exercises will fill out your hips and glutes to accentuate your figure. Upper body exercises will improve your posture and support your breasts which will make them look bigger. In other words, everything. You’re on estrogen. Have you seen cis women athletes? Exercise will feminize you.

\href{https://docs.google.com/document/d/1-NyE5EY5TTaRRMhk7HlTbKJ7HifjEsA4jlDO1qKQVl0/edit?tab=t.0}{This guide was shared with me} \textcolor{red}{(Warning: Google Docs link)} and looks to be a good starting place. I will note that there aren't particular exercises that feminize vs masculinize as bodies don't work like that, but you may wish to focus more on lower body exerices and flexibility more than the typical lifter.

\subsection{Can estrogen really cause height shrinkage?}

Yes. It is possible that it’s related to water content changes within tendons and ligaments, but it is not something that has been studied so the cause is fully speculation. Scientists: free study idea!

\subsection{Can estrogen really cause foot shrinkage?}

Yes. See above.

\subsection{Can estrogen really cause any other kinds of shrinkage?}\label{11-19}

Well, “use it or lose it” like they always say.

More seriously: without the presence of androgens such as testosterone regulating penile health via nocturnal erections and random erections, you are likely to experience far less erections without active effort. Without regular blood flow and oxidation via erections, the spongy tissue that composes the penis will atrophy over time via fibrosis. This leads to reduced firmness and possible pain during erections, potentially discouraging future erections and perpetuating the cycle of atrophy, which may cause issues during sex and/or for certain types of bottom surgery (namely: penile inversion vaginoplasty). It is not clear to what degree (if any) this is permanent, but the recommendations listed below in Question \ref{11-20} are more than sufficient for adequate mitigation of any issues. In simpler terms: it is a muscle to be exercised like any other.

\subsection*{Sexual Health}
\addcontentsline{toc}{subsection}{\textemdash{} Sexual Health}

\subsection{How do I improve erectile function on HRT?}\label{11-20}

Aside from using it regularly, ways to improve erectile function include: 1) Improving your fitness and physical health, particularly your cardiovascular ability; 2) consider medication like \textit{tadalafil} or \textit{sildenafil}; and 3) consider testosterone supplementation (see Section \ref{T} “TESTOSTERONE”). \textbf{Erectile dysfunction is not a guarantee on HRT and is entirely mitigable if you wish.}

For context on the two medications listed: \textit{tadalafil} is taken daily which improves the ability to have an erection at any time, whereas \textit{sildenafil} is essentially an erection on demand. They each have their own dosings and possible side effects not listed here.

If you would like to read a longer explanation for how erectile function works, \href{https://stainedglasswoman.substack.com/p/how-to-maintain-your-penis-function}{this Substack article} provides a good overview of the topic.

\subsection{How do I increase cum/pre-cum volume on HRT?}

Don’t be embarrassed, it’s a common question. Sunflower lecithin and pygeum increase both of those, in addition to testosterone supplementation as mentioned above. It seems to also make a difference for vaginal wetness and arousal for those who have had bottom surgery, but data and anecdotes are limited so it’s hard to say. Otherwise just be sure you drink enough water and have your nutrition in check.

\subsection{Can I lactate on HRT?}

Yes. Domperidone, fenugreek, sunflower lecithin, ample estrogen, and ample progesterone. Get a pump. Knock yourself out.

It should be noted that domperidone has side effects and risks associated with it, and that ability to lactate does not affect breast development. Newman-Goldfarb protocols would be what you want to look into.

\subsection{Can HRT change your senses and your perceptions, i.e. smell?}

You very likely were dissociated and depressed for years prior to starting HRT. The world is more vibrant now because you are no longer dissociating 24/7. The wonders of modern medicine!

It can, however, directly change your eye prescription. That can definitely happen.

\subsection{Can HRT change your sexuality?}

Similar to being dissociated as with above, HRT often incurs a lot more openness and acceptance with yourself which can cause a shift in how your sexuality presents itself. It is largely a semantics argument as to whether that is chemical or behavioral. A matter of perspective.

\subsection{Should I be on PrEP?}\label{11-25}

\textbf{Yes.} \href{https://en.wikipedia.org/wiki/Pre-exposure_prophylaxis_for_HIV_prevention}{\textit{Pre-exposure prophylaxis for HIV prevention} (PrEP)} is a category of antiviral drugs for the purpose of preventing HIV/AIDS. This is not directly related to HRT, but it is common for trans women to be at elevated risk of HIV/AIDs. Given the history of the AIDS pandemic, PrEP is a miracle of modern medicine that should interest you. \textbf{Note: There is no effect of any PrEP drug on HRT so you are encouraged to be on PrEP.}

If you are sexually active, you should strongly consider being on PrEP. Even if you are not sexually active, trans women are at a significantly higher risk of sexual violence, so you should still strongly consider being on PrEP. In most places you are likely to be prescribed \textit{Truvada} as a once-daily pill, although if you experience nausea as a side effect you can likely swap to \textit{Descovy} with no change in effectiveness. In the US, insurance commonly does not cover \textit{Descovy} unless you have tried or say you have tried \textit{Truvada}. The novel drug \textit{lenacapavir} as a twice-yearly injection is expected to make PrEP significantly easier as access becomes more widely available, if current options are prohibitive for you.

\subsection*{Medical Malpractice}
\addcontentsline{toc}{subsection}{\textemdash{} Medical Malpractice}

\subsection{I heard that injections are actually less stable because you do them less frequently. Is that true?}

Only if you follow the dipshit WPATH SOC 8 guidelines that list a recommended regimen of \textit{estradiol valerate} or \textit{estradiol cypionate} in the range of 5-30mg every two weeks which, to be abundantly clear, you absolutely should never do in a million years. “Do no harm”, my ass.

\subsection{But my doctor said-?}

The average doctor has essentially no training in anything related to trans healthcare, and \href{https://www.endocrine.org/news-and-advocacy/news-room/2017/endocrinologists-want-training-in-transgender-care}{4/5 endocrinologists have never had any formal training in trans healthcare}. It is most likely that you are their first trans patient and that they are inexperienced in the practical elements of managing a trans patient. Even among doctors who care a lot, they are often limited by conservative standards of care that they are forced to follow which do not always align with the care best for you. See above.

Please also be aware of “trans broken arm syndrome”, aka the tendency of doctors to blame everything on HRT. If your arm is broken, it's probably not “because of those hormones”!

And I should put this as a separate question but I don't want to break the formatting: in line with medical malpractice, there is no situation in which it is reasonable for a doctor to request to see or feel your breasts to “monitor growth” or for any other reason. It is far less common these days, thankfully, but it is sexual assault and completely unacceptable.

\subsection{My doctor won’t prescribe me injections. What do I do?}

Attempt to convince them, replace them, or seek DIY sources. Do not let a gatekeeping medical establishment prevent you from receiving the appropriate care that you deserve. \textbf{The most crucial aspect of interfacing with the medical system while trans is that you have to advocate for yourself.} This is compounded with disability, ethnicity, and other afflictions that scare doctors like womanhood.

\subsection{How does HRT for menopausal cis women relate to HRT for trans women?}\label{11-29}

While we generally have different goals and crucially have very different dosage requirements, there is an immense amount of overlap in experience for trans women and menopausal cis women. Medical misogyny in the form of incompetence, dismissiveness, antagonism, and/or misinformation is something that we unfortunately both experience. It is for this reason that it is paramount to build solidarity on this front. To give an example of what I mean, \href{https://www.youtube.com/watch?v=W0XW6av2wLQ}{the first 30-40 minutes of this interview} will likely sound extremely familiar to you if you would like to raise your blood pressure. The interviewee herself notes the connection too! The WHI ruined the lives of countless women.

\subsection*{Intersexuality and Comorbidities}
\addcontentsline{toc}{subsection}{\textemdash{} Intersexuality and Comorbidities}

\subsection{What’s up with Ehlers-Danlos Syndrome?}

This connective tissue disorder doesn’t actually relate to HRT but a lot of trans people have it so congrats in case this is how you learned that you do too. Aside from general cardiovascular long term concerns to maybe look into, keep up with strength training so that your joints work. Look into that elsewhere though. See Question \ref{11-16}.

\subsection{What kind of intersex things should I keep in mind?}

Throughout this guide, I have mentioned intersex conditions vaguely. Below is a short list of things that might be useful for you to know in your travels for yourself or for a friend.

\subsection{What’s up with Klinefelter Syndrome?}

This is a relatively (considering chromosomal mutations) common intersex-related condition that some trans women might not realize that they have as the two can overlap. It generally presents as low testosterone at the start of puberty. Good for you to know the name, just in case.

\subsection{What’s up with Persistent Müllerian Duct Syndrome (PMDS)?}

Another “I’m putting this here because this might be the first time you’ve even heard of the term” intersex-related condition that can affect some trans women, however few that may be since we don’t have numbers. The possible presence of an underdeveloped uterus leads to some possible complications and oddities. You probably extra want to have progesterone to avoid uterine cancer risks.

\subsection{What's up with ovotesticular syndrome?}

This intersex condition in particular can cause early level fluctuations which made lead to confusing test results due to the presence of both ovarian and testicular tissues, either separate or combined in an \textit{ovotestis}. This presents in many different ways which HRT can interact with as you begin suppressing \textit{luteinizing hormone} (LH). A uterus may or may not be present, multiple sets of gonads could be present, and/or it could look outwardly typical.

\subsection{What’s the difference between intestinal cramps and uterine cramps?}\label{11-35}

These are commonly misattributed in early transition as a symptom of intersex conditions. Intestinal cramps are widespread and diffuse across your abdomen, whereas uterine cramps are highly concentrated in a location somewhere below your belly button and tend to be sharp stabs/contractions in rapid succession. Like the inside of your body is used as a stress ball. Very different!

\subsection{What about other intersex conditions?}

I have listed a few notable ones, but there are far more expressions and ways of testing them that go far beyond the scope of this guide. Anecdotally, prevalence is higher than average among trans people so basic familiarity with this is useful.

\subsection*{Oddball Questions}
\addcontentsline{toc}{subsection}{\textemdash{} Oddball Questions}

\subsection{Many DIY sources only take crypto. Is that required? How does that work?}

There are other guides that cover this in better depth than I can on how to use crypto safely, including some vendors who have their own guides. But yes, crypto is often required for a lot of reasons. “Crypto” means a lot of things, but using it as a currency was the original point after all. It’s mostly just a pain in the ass. Monero (XMR) is good.

\subsection{What about Selective Estrogen Receptor Modulator (SERM) drugs for nonbinary regimens?}

Some people use SERMs as a part of a transition that is not looking to feminize as much for a more androgynous look, but it’s pretty much entirely uncharted waters thus why their mention is otherwise absent from this guide. You’re on your own if that’s something you want to explore, so please be safe. I don’t personally rate them very highly as I have not seen much to suggest that they work well for how people usually think or want them to work, at least not without a lot more caveats, but obviously there are people who like them. It's just not something I feel comfortable giving recommendations for.

The various proposed nonbinary regimens are often highly individualized because they are specific to what a persons' particular goals are. All HRT should be individualized to a degree, but there is often more variation in desired outcomes when people ask about androgyny. Hormonally, it is nontrivial. Everything stated in this guide should be treated solely as a starting place if you are wanting to experiment with something more complicated, but do remember that there is much more to achieving transition goals than just hormones alone.

\subsection{Are things like “herbal HRT” or “phytoestrogens” legitimate?}

\textbf{No.} If someone is telling you they have “herbal HRT”, they are telling you they have snake oil. The only thing that is going to feminize you is estrogen, not plant estrogens. No amount of “natural” products are a replacement for estrogen itself. This isn’t a common scam and you probably already know, but just in case you run into it, now you know for sure. If it smells like bullshit, it’s probably bullshit. Unless we’re talking about bug steroids in which case yeah those are actually cool. Won’t feminize you though.

\subsection{Is the Reddit Doctor that people constantly talk about Good?}

No.

\subsection{I hear DIY estrogen is made in a bathtub. Is that true?}

No. I honestly have no idea where or why this joke started that people now take seriously, but there’s no step in any process where a bathtub would even be considered. Don’t believe everything you read online. I don’t even know what you could even theoretically do with a bathtub, unless you think estrogen vials are full of the bathwater of trans women. I don’t know why you would think that though. It’s obviously cum.

\subsection{How does HRT affect fertility?}\label{11-42}

It is important to understand that this is extremely understudied so exact figures cannot be stated, and given the seriousness of pregnancy, I urge you to practice safe sex and lean on the side of caution where possible. HRT itself can, and likely will, make you infertile eventually, but only through full suppression of \href{https://en.wikipedia.org/wiki/Hypothalamic-pituitary-gonadal_axis}{the HPG axis} (see Question \ref{2-3}) over a long time span. In other words, if you haven't had bottom surgery of any kind and you are on an HRT regimen that is less capable of HPG axis suppression (such as pills), then this is more of a consideration.

\textbf{If the HPG axis is not suppressed then it is fully possible to impregnate someone}, and the timeline for sperm maturation is long enough that this is true even after the HPG axis has been initially suppressed for \textbf{multiple months}. Please take this very seriously. Full HPG axis suppression for at minimum six months, perhaps closer to a year out of an abundance of caution, is recommended.

\subsection{Is infertility from HRT reversible?}\label{11-43}

It is theoretically possible to reverse HRT-induced infertility, assuming you weren't already infertile prior to HRT (a large assumption!), but there are not many documented cases of this so the full efficacy of fertility restoration after long-term HRT is unknown. The process would involve restarting the HPG axis with a variety of medications along with entirely stopping HRT, which would in essence require a hormonal detransition for likely six months at minimum, and even then sperm quality is not certain or guaranteed. It is not something that should be planned for, to say the least, so planning around it would be wise. A sperm bank would be recommended before or early in HRT, financially permitting, if potential biological children are a priority and if a future relationship where that is possible/desired is likely.



\section{CREATINE}

\subsection{What is creatine?}

Creatine is an organic compound in your muscles and in your brain. It recycles ADP into ATP which is important for energy production in your body, especially initial high burst applications before other energy systems take over.

\subsection{Isn’t it like a steroid or something that bodybuilders use?}

No. Bodybuilders and athletes like it because having more energy means more activity before getting tired. They aren’t the only ones who use it since it is basically the \#1 supplement in terms of things that are actually useful and are actually researched.

\subsection{How is creatine related to HRT?}

\textbf{It isn’t!} But it’s something I yell about because I think it’s good and I am tired of repeating myself because people keep asking and you’re reading this anyway, aren’t you? I love a captive audience. My standup routine is at the bottom.

\subsection{Okay well why should I take creatine then?}

What a great question! It’s good for your brain and your muscles. Creatine is often found in relatively low concentrations for many people depending on their diet, especially people who don’t eat meat. There is compelling research about various chronic fatigue and post-viral conditions (long COVID in particular) being related to depleted creatine reserves in the brain, so some people find cognitive benefits from supplementing it. It isn’t magic but it is dirt cheap so it is worth trying in my opinion.

\subsection{What are the forms?}

Just \textit{creatine monohydrate} powder is what you want. The pills tend to be low dosage and are up charging you anyway, while gummies often destroy the creatine in the creation of the gummy. A lot of brands include creatine in some sort of mix but the pure stuff is usually cheaper.

\subsection{How do I take it then?}

The general recommendation is 5-10g daily dissolved in some sort of liquid. It dissolves best in things that aren’t just water. It’s mostly flavorless, so just throw a scoop or two in your coffee or a smoothie and call it a day. It can be a little chalky or gritty depending on the quantity and the fluid.

\subsection{Does it matter when I take it?}

Not really. It doesn’t have an immediate effect like that which is why it’s silly that it’s microdosed in pre-workout mixes. Take it whenever it’s convenient for you.

\subsection{How does it work then?}

It builds up in your body to a maximum level of saturation over a week or two. Then you just maintain that and reap the rewards (of maybe feeling better).

\subsection{Do I have to do a “loading” phase of taking more at first?}

Probably not. Unless you’re on some sort of intense training time crunch or something, this probably doesn’t matter at all. Just take whatever is convenient with some regularity.

\subsection{What are the side effects?}

Slight weight gain may be possible because of increased water weight in your muscles (which to be clear is Good, so don't be alarmed). If you don’t take it with water, or if you take too much at once, you might get a tummy ache. Ouchie.

\subsection{Who shouldn’t take it?}

People with kidney issues. Not because it causes them, but because creatinine (different spelling! Creatine becomes creatinine) is used as a marker in lab tests for a number of kidney issues and supplementing might give a false positive. Just keep it in mind.

\subsection{Do you have any brand recommendations?}

No. It shouldn’t really matter. Just get whatever seems reputable and is a reasonable price. I’d give a recommendation for the one I like but when I asked the brand for affiliate link they turned me down, so their loss! No free clout.

\subsection{You seriously put creatine into this document, huh?}

Yeah it’s pretty funny. It’s not my fault that I joked about it and people told me it legitimately helped them because now I feel obligated to keep talking about it!!!



\section{CLOSING REMARKS}

If any of the following are true:

\begin{itemize}
\item you are still mad at me despite the disclaimer;

\item you spotted an issue or typo;

\item you have a clarifying question that should be put into the text;

\item you have an objection that hopefully isn’t an Uhm Ackshually;

\item you wish to sing my praises;

\item you wish to pledge fealty;

\item you wish to send tithes my way;
\end{itemize}

Then please feel free to contact me and I’ll see what we can do. Bluesky is the easiest contact point, and you can DM me for my Signal. Otherwise, thank you for reading and I hope it helps.

\textbf{If you would like to donate to support this project,} \href{https://cash.app/Katitties}{CashApp}, \href{https://ko-fi.com/katitties}{Ko-Fi}, and \href{https://account.venmo.com/u/katitties}{Venmo} all work. I appreciate it!

And lastly: \textbf{The most important thing that you can do as a trans person is to live.} For as much as this document is a manual, it is in equal measure a message to you as a trans person that your existence is a gift upon the world, your presence is a blessing on those around you, and that you deserve to be treated with respect. Even if you do nothing else, your life is a feat worth praising. Thank you.



\section*{FRIENDS OF PGHRT}\label{FOPGHRT}
\addcontentsline{toc}{section}{FRIENDS OF PGHRT}

Across this document is a scattering of links to other guides and resources. Below is a consolidation of them which will also include more links to external resources as time goes on, ideally by other trans people. For the privacy minded or noided, note that some of these are Google Docs links.

\begin{enumerate}
  \item \href{https://startwith4mgestradiolenanthateweeklyandtestatonetothreemonths.com/}{SW4EEWATAOTTM} - TL;DR for PGHRT
  \item \href{https://hrtcafe.net/}{HRT Cafe} - HRT Resource Aggregator
  \item \href{https://transfemscience.org/}{Transfeminine Science} - Informational resource for trans medical literature
  \item \href{https://estrannai.se}{Estrannai.se} - Estradiol Pharmacokinetics Playground
  \item \href{https://globoho.moe/}{Globoho.moe} - Thailand Orchiectomy Medical Tourism Travel Guide
  \item Julia's FUE Guide - COMING SOON, I'M BULLYING HER TO WRITE FASTER
  \item \href{https://docs.google.com/document/d/1-NyE5EY5TTaRRMhk7HlTbKJ7HifjEsA4jlDO1qKQVl0/edit?tab=t.0}{Sky's Feminine Figure Beginner Program} - An exercise regimen geared towards trans fems
  \item \href{https://docs.google.com/document/d/114sztSw1aVWM2pXLDl9NrHklyvewz3EmFiHiisjM71k/edit?tab=t.0}{Sky's Diet 101} - A guide towards adjusting weight in a healthy way
  \item \href{https://stainedglasswoman.substack.com/p/how-to-maintain-your-penis-function}{How to Maintain Erectile Function on HRT} - A longer form explanation on the “use it or lose it” phenomenon
  \item \href{https://mega.nz/folder/cfwSmBbY#ify6kaueB8SdDD1v12t2Dw}{Biohax Guide} - Trans Masc DIY Guide
\end{enumerate}

\section*{ABOUT THE AUTHOR}
\addcontentsline{toc}{section}{ABOUT THE AUTHOR}

Katie Tightpussy is an award-winning author and professional trans woman with nearly a decade of experience in the field of transgender. Her accomplishments include transiferating her sex through the novel technique of cross-sex hormone injections, being physically unable to shut up, and utilizing a very fortunate set of hyperfixations as they relate to transbobulation of the humors. She spends her days in the idyllic rural countryside of Los Angeles scheming of new ways to achieve world domination and enjoys riding her bicycle. Media inquiries can reach her agent at \href{http://katietightpussy.com}{katietightpussy.com}.



\section*{DISCLOSURES}
\addcontentsline{toc}{section}{DISCLOSURES}

No robot girls were harmed in the making of this document, including any usage of generative large language models. The author does not endorse any reproduction without attribution nor scraping of this work. Leave those poor robot girls alone.

The author declares an attraction towards women and acknowledges a potential conflict of interest for the existence of more beautiful trans women in the world.



\section*{ACKNOWLEDGEMENTS}
\addcontentsline{toc}{section}{ACKNOWLEDGEMENTS}

Though the text is primarily my voice, this document would not be even half as good without the contributions, feedback, and suggestions from others involved at every step along the way. A good reminder as ever that transition is not something best done alone.

Many thanks to Q, R, RM, and S in alphabetical order for close review and generally being fun nerds to talk to; love y’all. Special thanks to CB and J for close review that also inspired some very good bits. Thanks to KG for additional intersex information. Thanks to w [sic] for additional injection resources. Thanks to BIR collectively for a plethora of crucial nerd nitpicks. Appreciation for general review from c [sic], JTP, K, S, and V. Thanks to E and S for web dev assistance. Thanks to everyone on Bluesky who encouraged me to write this up in the first place, and everyone over the years sharing knowledge. And of course: much appreciation to all HRT nerds, even when we disagree, since we’re all trying to do the best for our community where we’ve otherwise been let down. Keep up the good work everyone.

Shout out to my IB Chemistry HL teacher many years ago who quite reasonably doubted my studiousness even though I’m now putting much of that knowledge to use for the art of transsexuality; go figure.



\section*{CHANGELOG}
\addcontentsline{toc}{section}{CHANGELOG}

\noindent \href{https://github.com/Juicysteak117/pghrt/}{Source code available here on GitHub.}

\noindent Full Compilation Datetime: \DTMnow

\noindent(There aren't LaTeXML bindings for \texttt{datetime2}, \texttt{hanging}, or \texttt{hyphenat}, so the formatting is slightly ugly. If you'd really like to help me out, please write those bindings!!!)

\noindent 2025-08-20: Initial release. 15.9k words.

\noindent 2025-08-20: A lot of typos and minor verbiage tweaks. Added Question \ref{8-18}.

\noindent 2025-08-21: Typos grow on trees. Added Question \ref{5-27}.

\noindent 2025-08-21: More tweaks. Opted to remove “WHY PROG” from Question \ref{8-17}. 17.0k words.

\noindent 2025-08-22: Nitpicks, clarifications, and typos. 17.2k words.

\noindent 2025-08-24: A few more twinks sorry tweaks. 17.3k words.

\noindent 2025-08-27: How long until remaining typos are embarrasing? 17.3k words.

\noindent 2025-08-28: Reduced ambiguity in a few areas. 17.4k words.

\noindent 2025-08-29: Additional clarity for frequencies in Section \ref{td}. 17.5k words.

\noindent 2025-09-01: Sisyphus boulder meme captioned fixing typos dot png. 17.5k words.

\noindent 2025-09-07: Added donation links per request. That's very kind. 17.5k words.

\noindent 2025-09-07: Few more tweaks. Clarified an additional common progestin. 17.6k words.

\noindent 2025-09-19: Added Question \ref{4-16} plus tweaks. 17.7k words.

\noindent 2025-09-23: A wide variety of clarifications up and down the line. 18.1k words.

\noindent 2025-09-24: Added an important note about surgery to Question \ref{11-1}. 18.3k words.

\noindent 2025-09-24: “Katie my doctor told me-” It never ends. 18.5k words.

\noindent 2025-09-26: Another pass of clarification edits. Yes I should have a git diff. Sorry that I don't. I thought I'd be done by now anyway! 18.7k words.

\noindent 2025-09-30: Added some cross references for clarity. 18.7k words.

\noindent 2025-10-02: More cross references. Likely will do another pass. 18.8k words.

\noindent 2025-10-02: Added a big bold warning about recapping to Question \ref{5-13} because SOMEONE didn't watch the video smh. 18.9k words.

\noindent 2025-10-10: Added Question \ref{11-42} and Question \ref{11-43} per request. Honestly I just forgot about fertility being a thing lol. Also added the Friends Of PGHRT postword section. 19.5k words.

\noindent 2025-10-10: Added Gretchen's Version (.txt) and fixed formatting. 19.5k words.

\noindent 2025-10-11: Added permalinks to everything, yay! And finally made a git repo. Look at me being a big girl, wow. 19.5k words.

\noindent 2025-10-11: Added an external link to Question \ref{11-20}. 19.6k words.

\noindent 2025-10-14: Reduced the pithiness and expanded the usefulness of Question \ref{11-25} per repeat request because I \textit{guess} there is public health utility in doing so. 19.8k words.

\noindent 2025-10-17: Dark mode! And a font toggle off. Plus a few links. 19.8k words.

\noindent 2025-10-25: A variety of micro tweaks and the inclusion of Question \ref{1-12}. 20.2k words

\noindent 2026-01-10: Added a clarification line under Question \ref{5-25} because that should be written somewhere. 20.3k words

\noindent 2026-01-24: Slight retweak to Question \ref{5-26} to better clarify how dead space works. 20.4k words.

\noindent 2026-01-29: OKAY FINE. Elaborated substantially on Question \ref{11-19} and Question \ref{11-20}. 20.6k words.

\noindent 2026-02-04: Slight common sense addendum to Question \ref{5-5}, Question \ref{5-9}, and Question \ref{5-17}. 20.8k words.

\end{document}